\documentclass[11pt,letterpaper]{article}
\usepackage[T1]{fontenc}
\usepackage[utf8]{inputenc}
\usepackage{lmodern,microtype}
\usepackage[margin=.92in,headheight=15pt]{geometry}
\usepackage{amsmath,amssymb,amsthm,mathtools}
\usepackage{booktabs,longtable,tabularx,array,enumitem}
\usepackage{xcolor,fancyhdr,listings,xurl,needspace}
\usepackage[colorlinks=true,linkcolor=blue!40!black,urlcolor=blue!45!black,citecolor=blue!40!black]{hyperref}
\usepackage{bookmark}
\definecolor{ink}{RGB}{25,45,65}
\definecolor{pale}{RGB}{246,248,251}
\setlength{\parskip}{.3em}
\setlength{\emergencystretch}{3em}
\setlist{itemsep=.25em,topsep=.45em,leftmargin=1.7em}
\newtheorem{theorem}{Theorem}[section]
\newtheorem{lemma}[theorem]{Lemma}
\newtheorem{proposition}[theorem]{Proposition}
\newtheorem{corollary}[theorem]{Corollary}
\theoremstyle{definition}\newtheorem{definition}[theorem]{Definition}
\theoremstyle{remark}\newtheorem{remark}[theorem]{Remark}
\newcommand{\ipra}{\mathsf{iPRA}}
\newcommand{\D}{\mathsf D}
\newcommand{\Q}{\mathsf Q}
\newcommand{\FV}{\operatorname{FV}}
\newcommand{\file}[1]{\expandafter\nolinkurl\expandafter{\detokenize{#1}}}
\lstset{basicstyle=\ttfamily\footnotesize,columns=fullflexible,keepspaces=true,
  breaklines=true,showstringspaces=false,frame=single,rulecolor=\color{black!15},
  backgroundcolor=\color{pale},aboveskip=8pt,belowskip=8pt}
\pagestyle{fancy}\fancyhf{}
\fancyhead[L]{\small A compact iPRA kernel and a symmetry obstruction}
\fancyhead[R]{\small 10 November 2026}
\fancyfoot[C]{\thepage}
\hypersetup{pdftitle={A Smaller Complete iPRA Checker and an Obstruction to Positive Axiomatizations},
 pdfauthor={Research prepared for Vladimir Reshetnikov},
 pdfsubject={Intuitionistic primitive recursive arithmetic; full bi-interpretability; reproducible checker size; conjunctive induction}}
\title{\color{ink}\bfseries A Smaller Complete iPRA Checker\\
and an Obstruction to Positive Axiomatizations\\[8pt]
\large Explicit rules, both interpretation composites,\\
a pinned source audit, and a quantified symmetry theorem}
\author{Research prepared for Vladimir Reshetnikov}
\date{10 November 2026}
\begin{document}
\maketitle
\begin{abstract}
We give a complete proof presentation of full first-order intuitionistic
primitive recursive arithmetic and an independently executable checker
with 612 normalized Python bytes, 266 significant tokens, and 310 syntax-tree
nodes. Maximum-compression gzip occupies 360 bytes and is reported only as
information. The current supplied corpus has separate minima of 616 bytes,
268 tokens, and 300 nodes. Thus the selected checker improves source size
and tokens, but not the syntax-tree minimum. Its ordinary list check does
not rely on divergent malformed-string processing. The construction
combines the existing arity representation with the existing predecessor
computation interface; it is not a new numerical or logical axiom basis.

The new mathematical contribution is a broader obstruction to positive
axiomatizations. Eventually constant powers have exact coordinatewise
truth for formulas built from atoms using conjunction and both quantifiers,
with arbitrary finite quantifier alternation. They satisfy induction for
all these formulas, all standard-true implications between them, and the
corresponding negative constraints. Nevertheless, after any finite set of
parameters is named, they retain a nontrivial involution. This excludes
ordinary first-order bi-interpretability with iPRA, even using additional
sorts, arbitrary formula-defined domains and graphs, finite tuples, and
definable quotient equality. The result extends an obstruction in the
supplied drafts and is not a source-size lower bound.

The archive supplies the complete primitive rule inventory, attributed
arithmetic and logical bridge proofs, executable proof certificates,
negative controls, an exhaustive source census, and reproducible measurements.
The broader request for a significant new simplest axiomatization remains
unresolved. No absolute minimality or proof-assistant verification is claimed.
\end{abstract}

\begin{center}
\fcolorbox{ink}{pale}{\begin{minipage}{.91\linewidth}
\textbf{Precise outcome.} A smaller complete implementation and a stronger
non-bi-interpretability theorem are established. The positive arithmetic
foundation is inherited. A four-byte implementation gain should not be
mistaken for the requested foundational breakthrough.
\end{minipage}}
\end{center}
\clearpage
\setcounter{tocdepth}{1}\tableofcontents
\clearpage


% Begin included source: theory/result.tex
\section{The result, the comparison, and the unresolved requirement}
\label{sec:result}

The intended target is full first-order \emph{intuitionistic} primitive
recursive arithmetic, not just its equational fragment or classical PRA.
The familiar numerical primitives and quantifier-free induction can be
presented with either classical or intuitionistic first-order logic;
that distinction matters for a claim of bi-interpretability
\cite{avigad,moschovakis}. We fix the intuitionistic theory throughout.
The full logical interface will be derived from an explicitly listed
native calculus rather than left as an unmeasured background theorem prover.

\subsection{A pinned and exhaustive implementation baseline}

The comparison uses commit
\begin{center}\small
\file{fb2dbd2a33c6b132671a54b029d0bf3cd98c633d}.
\end{center}
There are thirteen ZIP archives directly in \file{docs/drafts} at this
revision. Recursively including embedded historical archives gives 28
distinct archive byte strings and 370 distinct Python source byte strings.
The inventory is SHA-256 deduplicated and keeps every original archive
and member path. Small build scripts, test runners, and producer programs
are not complete proof checkers. Complete historical controls inside
\file{reference/} or \file{ancillary/} do count.

Every source at or below at least one final coordinate minimum was
individually classified. There are 31 such sources: three complete
checkers and 28 supporting programs. The current minimum byte count is
\emph{616}, not the older 672 quoted in several supplied reports. The
current minimum token and AST counts are 268 and 300. The complete
census and classifications are in \file{evidence/}.

\begin{center}\small
\begin{tabular}{lrrrr}
\toprule
Complete standalone program & Bytes & Tokens & AST & Gzip 9\\
\midrule
Current byte leader &616&269&311&366\\
Current AST leader &649&268&300&367\\
Current ordinary-list predecessor checker &629&269&307&373\\
Current matched old-flat control &628&269&306&374\\
\midrule
Selected arity/predecessor composition &\textbf{612}&\textbf{266}&310&360\\
Ancillary shorter-byte composition &597&267&315&356\\
Ancillary smaller-AST composition &639&267&298&360\\
\bottomrule
\end{tabular}
\end{center}

Bytes means the size after removing genuine docstrings and comments,
applying \texttt{ast.unparse}, and adding one final newline. Significant
tokens exclude layout and comment tokens. AST counts include every node
of the parsed module. Gzip is measured on the exact delivered raw source
with \texttt{gzip.compress(data, compresslevel=9, mtime=0)}; the complete
gzip member, including header and trailer, is counted. Compression is
not an optimization objective. The runtime, source hashes, raw sizes,
and normalized-source gzip sizes are recorded with the measurements.

These are whole-program measurements. They include the JSON input,
traversal and scope checks, arithmetic supply, logical rules, recursive
proof checking, and success output. The Python runtime and standard JSON
library are common infrastructure. No producer, theorem library,
arithmetical evaluator, or precomputed rule table is imported by the
selected program. Its separately measured readable expansion is also
executable.

The selected program improves the current byte coordinate by
$616-612=4$ bytes, approximately $0.65\%$, and the token coordinate by
two tokens. It has ten more AST nodes than the separate AST minimum.
The two ancillary programs show additional tradeoffs; they do not make
the selected program a simultaneous record in all coordinates.
The smaller-AST opener inherits a qualification concerning malformed
strings and CPython's recursion limit, so the ordinary-list program is
the main artifact.

\subsection{What is new and what is inherited}

The arithmetic basis $0,S,I$ and the native Eq/$\Q$/$\D$ inference
rules are inherited. The arity codec is inherited from~\cite{arity},
the predecessor packet from~\cite{uniform}, and the exit-status
rejection convention already appears in~\cite{uniform}. We prove that
these mechanisms compose, including an exact simulation of arithmetic
packets and a direct justification of zero-witness introduction.
This supplies a useful smaller complete checkpoint, but does not
establish a new foundational axiom basis.

The mathematical extension is the theorem of
Section~\ref{sec:pp-obstruction}. The earlier eventually constant model
argument covers equations and finite conjunctions. Exact transfer in
the same model persists through both quantifiers, even with arbitrary
finite alternation. Consequently even a much richer class of positive
axioms and induction schemas leaves an obstruction to
bi-interpretability. The interpretation may use a different domain,
finite tuples, multiple sorts, arbitrary definable graphs, and quotient
equality. The argument does not rely on term-definability of witnesses.

\subsection{Why the word ``lower bound'' needs qualification}

A minimum among the submitted implementations is an observed envelope.
It is an upper bound on the unknown optimum, not a lower-bound theorem
for every possible axiomatization. Source size also depends on a fixed
programming language, runtime interface, normalization convention, and
certificate encoding. We preserve the corpus's conventions so that the
comparison is reproducible, and do not infer absolute simplicity from it.

The requested significant new simplest axiomatization has not been
established. The numerical and logical presentation below is exact and
bi-interpretable with iPRA, but its basis is already present in the
drafts. The size improvement and new obstruction are stated separately
so that neither is used to overstate the other.

% End included source


% Begin included source: theory/presentation.tex
\section{A complete native presentation of full intuitionistic PRA}
\label{sec:presentation}

This section specifies the mathematical theory checked by the program.
It is the inherited Choice-DI scalar-iteration presentation
\cite{uniform,arity}. Its rules, numerical operations, and interpretation
proofs are retained explicitly. The new executable is a representation
of this theory, rather than a different arithmetical theory whose
equivalence is only conjectural.

\subsection{Numerical syntax and the two computation schemas}

There is one numerical sort, with finite terms
\[
 t ::= x\mid 0\mid St\mid I(n,a;u.h).
\]
Only the update body $h$ binds $u$. The counter $n$ and initializer $a$
do not bind it. Other variables in $h$ are passive numerical parameters
and remain fixed during iteration. In every retained context, the
following are axiom schemas:
\begin{align}
 I(0,a;u.h)&=a, &&\text{(I0)}\label{eq:new-I0}\\
 I(Sn,a;u.h)&=h[I(n,a;u.h)/u]. &&\text{(IS)}\label{eq:new-IS}
\end{align}
The update sees only its previous value, not the old counter.

The binder is a finite-description notation with an ordinary first-order
elaboration. Beginning with $0,S$, each earlier-stage term
$h(u,\vec p)$ admits a new symbol $i_h(n,a,\vec p)$ with the two
displayed defining equations. Each expression and derivation involves
only finitely many stages. This introduces neither a function sort nor
an object variable ranging over programs. In particular, it does not
postulate a total universal evaluator for all primitive-recursive
descriptions.

\subsection{Formula syntax and all primitive logical rules}

Formulas are finite expressions
\[
 A ::= (s=t)\mid \Q_y A\mid\D_{x,y}A.
\]
The variables $x,y$ bound by $\D$ are distinct. Contexts $\Gamma$ are
finite reusable lists of formulas. Substitution always avoids capture;
local bound variables are renamed as needed. The designated
contradiction is the equation $0=S0$. The intended readings, which
will be proved exact rather than assumed as extra inference rules, are
\[
 \Q_y A(y)\equiv A(0)\to\forall y A(y),\qquad
 \D_{x,y}A(x,y)\equiv
 \exists x\bigl(A(x,0)\to\forall y A(x,y)\bigr).
\]

Together with I0 and IS, the following list is complete:
\[
 \frac{}{\Gamma\vdash A}\;(A\in\Gamma)
 \quad\text{(Hyp)},
\]
\[
 \frac{\Gamma\vdash s=t\qquad\Gamma\vdash F(s)}
      {\Gamma\vdash F(t)}
 \quad\text{(Tr)},
\]
\[
 \frac{\Gamma\vdash\Q_yA(y)\qquad\Gamma\vdash A(0)}
      {\Gamma\vdash A(t)}
 \quad\text{(QE)},
\]
\[
 \frac{\Gamma\vdash\D_{x,y}A(x,y)\qquad
       \Gamma,\Q_yA(b,y)\vdash F}
      {\Gamma\vdash F}
 \quad\text{(DE)},
\]
\[
 \frac{\Gamma,A(t,0)\vdash A(t,Sz)}
      {\Gamma\vdash\D_{x,y}A(x,y)}
 \quad\text{(DI-case)},
\]
\[
 \frac{\Gamma,E(t,z)\vdash E(t,Sz)}
      {\Gamma\vdash\D_{x,y}E(x,y)}
 \quad\text{(DI-progress)}.
\]

In Tr, $F(v)$ is any fixed native formula template, including a
compound one. In DE, $b$ is fresh for $\Gamma,F$ and the fixed
parameters of the matrix. In both introductions, $z$ is fresh for
$\Gamma$, the chosen witness $t$, and the fixed matrix parameters.
The displayed local assumptions are discharged. In DI-progress,
$E(x,y)$ must be a \emph{literal equation}. An equation inside a
quantified formula does not qualify. DI-case is available for both
equations and compound formulas.

There are no additional primitive equality, successor, induction, or
ordinary logical axioms or rules. In particular, unrestricted progress
on compound formulas is not part of the theory. The program checks the
root-constructor restriction directly through its arity representation.

\subsection{First derived principles}

Weakening, exchange, contraction, renaming, numerical substitution,
and ordinary-assumption cut are admissible by induction on finite
derivations. Rename eigenparameters before inserting a proof or term.
Numerical substitution leaves an induction matrix's root constructor
unchanged, so the equation-only restriction survives these transformations.

Reflexivity is derived. Let $B(p)=I(0,p;v.0)$, so I0 proves $B(p)=p$.
Use that equation as both premises of Tr with template $F(w):=(w=p)$.
The result is $p=p$. Transport then supplies symmetry, transitivity,
and congruence for every fixed numerical expression, as well as transport
inside arbitrary native formulas.

The auxiliary closure $\Q_y A$ is interderivable with
$\D_{x,y}A$ when $x$ is unused. From $\Q_y A$, DI-case specializes
the closure under its base assumption. Conversely, DE opens the paired
formula to exactly $\Q_y A$, and Hyp supplies the continuation. Thus
the public use of $\Q$ introduces no independent logical strength.

With an unused existential coordinate, DI-progress, DE, and QE derive
the contextual equation-induction rule
\[
 \frac{\Gamma\vdash E(0)\qquad\Gamma,E(z)\vdash E(Sz)}
      {\Gamma\vdash E(t)},
 \qquad z\text{ fresh}.
 \tag{EqInd}\label{eq:new-eqind}
\]
The construction first proves $\D_{x,y}E(y)$ by progress, opens its
closure by DE, and applies QE to the separately proved base. The
same unused-coordinate construction with DI-case gives numerical
zero/successor cases for an arbitrary formula:
\[
 \frac{\Gamma\vdash A(0)\qquad\Gamma\vdash A(Sz)}
      {\Gamma\vdash A(t)},
 \qquad z\text{ fresh}.
 \tag{Cases}\label{eq:new-cases}
\]
Only EqInd permits the current instance as an induction assumption.
These two rules must not be conflated.

\subsection{Recovering ordinary intuitionistic logic}

Set $1=S0$, $2=S1$, and define a numerical selector
\[
 C(v,a,b)=I(v,a;u.b).
\]
The computation schemas give $C(0,a,b)=a$ and $C(Sv,a,b)=b$.
These two endpoint identities are enough to construct a formula gate
by external recursion on each fixed formula:
\[
 G_v(s=t):=(C(v,0,s)=C(v,0,t)),\qquad
 G_v(\D_{x,y}A):=\D_{x,y}G_vA.
\]
First replace an auxiliary $\Q$ by its unused-witness $\D$ version.
Then $G_0A$ is provable and $G_{Sv}A$ is interderivable with $A$.
Padding by $\D_{u,y}G_yA$ preserves a formula and permits binder
depths to be aligned.

For equations put
\[
 M_v(s=t,r=w):=(C(v,s,r)=C(v,t,w)).
\]
For paired formulas align binders and recurse on their bodies. If just
one input is an equation, pad it first. The maximum remaining binder
depth decreases, so this is a finite external construction. Derived
binder congruence gives
\[
 M_0(A,B)\dashv\vdash A,\qquad M_{Sv}(A,B)\dashv\vdash B.
\]
Let $\delta(v)=C(v,1,0)$ and $p(z)=I(z,0;u.\delta(u))$.
Only the finite computations $p(1)=1$ and $p(2)=0$ are needed for
the following logical macros, with all unused coordinates fresh:
\begin{align*}
 A\to_*B&:=\D_{u,z}M_z(A,B),\\
 A\wedge_*B&:=\D_{u,z}G_zM_{p(z)}(A,B),\\
 \exists_*x\,A(x)&:=\D_{x,y}G_yA(x),\\
 A\vee_*B&:=\exists_*x\,M_x(A,B),\\
 \forall_*x\,A(x)&:=A(0)\wedge_*\D_{u,z}G_zA(z),\\
 \bot_*&:=(0=1).
\end{align*}

For implication introduction, the retained zero instance is the
antecedent and the successor instance is the consequent. For elimination,
open the closure and specialize at $1$. Conjunction introduction uses
the gate's provable zero case; its projections specialize the positive
gate at $1,2$. Existential introduction uses the proposed numerical
witness; elimination opens a fresh witness, proves the gate's zero
instance, and specializes at $1$. Disjunction uses numerical cases on
the selected witness. Universal elimination uses its explicit zero
conjunct or the successor closure; universal introduction uses a fresh
arbitrary parameter. Every ordinary introduction and elimination rule
is proved in the first appendix, including its eigenvariable conditions.

Numerical cases now prove
\[
 x=0\ \vee_*\ \exists_*y\,(x=Sy).
\]
When a target occurs in a retained context, one uses this equality-based
case split and Tr; one does not substitute into the retained context
without justification. Under $0=1$, transporting the selector endpoint
identities yields any numerical equation. External induction on formula
syntax using DI-case gives explosion for every native formula. This
recovers genuine intuitionistic falsity without adding excluded middle
for arbitrary formulas.

\subsection{Recovering all primitive-recursive arithmetic}

There is a substantial inherited arithmetic argument behind the small
primitive basis. The appendices give its finite terms and proofs, in
the following noncircular order:
\begin{enumerate}
\item Scalar iteration, cases, and literal-equation induction construct
      predecessor $P$, with $P(0)=0$ and $P(Sx)=x$.
\item A primitive-recursive characteristic of \emph{actual equality}
      is proved correct on the entire domain, not merely on standard
      numerals. Boolean characteristics of fixed quantifier-free
      formulas then reduce their induction to EqInd.
\item Quantifier-free induction proves the required pairing and
      projection identities. Pairing was not presupposed in the
      construction of predecessor.
\item Pairing compiles counter-aware primitive recursion into
      scalar iteration with a passive initializer.
\end{enumerate}

Here is the last compiler explicitly. Suppose a primitive-recursive
description has base $g(\vec p)$ and step $h(n,v,\vec p)$.
Let $\pi,\pi_0,\pi_1$ be the proved pairing and projections, and
assume the smaller descriptions have already been translated. Define
\begin{align*}
 V(n,\vec p)&=I\bigl(n,\pi(0,g^*(\vec p));\ u.
       \pi(S\pi_0(u),h^*(\pi_0(u),\pi_1(u),\vec p))\bigr),\\
 f^*(n,\vec p)&=\pi_1(V(n,\vec p)).
\end{align*}
Equation induction proves $\pi_0(V(n,\vec p))=n$. The projection
identities then prove the required base and step equations for $f^*$.
Composition and projections translate directly. The compilation is an
external recursion on a finite function description; it is not a single
object-language evaluator for all primitive-recursive programs.

Successor injectivity follows by applying predecessor. If $Sx=0$,
the selector distinguishes the two sides and derives $0=1$; hence
successor separation follows under the decoded falsity. Together with
coverage, full intuitionistic logic, the PR compiler, and quantifier-free
induction, this recovers the usual first-order axioms of iPRA.

\subsection{Both interpretations and both composites}

For a completely ordinary first-order comparison theory, let $T_I$
have the staged iterator symbols, intuitionistic first-order equality
logic, I0, IS, EqInd, and
\[
 x=0\vee\exists y(x=Sy),\qquad (0=S0)\to\bot.
\]
The last bridge is explicit. Positive computation equations alone have
a singleton ordinary model; naming $0=S0$ as native falsity is not a
reason to omit the bridge from $T_I$.

\begin{theorem}[Full-domain bi-interpretability]
\label{thm:new-equivalence}
The native calculus is a faithful presentation of $T_I$. It is
bi-interpretable and equiconsistent with full first-order iPRA. The
numerical domains are unchanged and equality is actual equality.
Both domain composites are internally isomorphic to the identity by
the relation $x=y$. Both formula composites are derivably equivalent
to the original formulas.
\end{theorem}
\begin{proof}
Interpret a fixed scalar iterator in iPRA by primitive recursion whose
step ignores the old counter. Preserve numerical equality and use
\[
 (\Q_yA)^\circ=A^\circ(0)\to\forall y A^\circ(y),\qquad
 (\D_{x,y}A)^\circ=
 \exists x\bigl(A^\circ(x,0)\to\forall y A^\circ(x,y)\bigr).
\]
Tr, QE, and DE are valid under these readings. DI-case is valid by
successor coverage; DI-progress is valid by equation induction, under
the temporarily assumed zero instance. I0 and IS are the primitive
recursion equations. Thus every native proof translates to an iPRA proof.

In the reverse direction, translate every PR function description by
the scalar compiler and every ordinary connective and quantifier by
the displayed macros. The derived rules and arithmetic bridge validate
each logical and nonlogical axiom and each inference of iPRA.

For numerical round trips, two scalar iterations with the same seed
and provably equal fixed updates are equal at every counter by EqInd.
The analogous uniqueness proof applies to each compiled primitive
recursion. External induction on a finite function description therefore
identifies its double translation with the original function, including
its passive parameters. This also proves coherence of the staged binder
notation. The domain is the whole numerical domain in each direction;
$x=y$ is consequently an internally total, bijective isomorphism of
the composite with the original structure.

For the nontrivial native formula composite, prove
\[
 \D_{x,y}A\dashv\vdash
 \exists_*x\bigl(A(x,0)\to_*\forall_*y A(x,y)\bigr).
\]
Forward, use DE at a fresh witness, assume the base, and specialize
the exposed closure at an arbitrary numerical argument. Reverse,
eliminate the translated existential, use its implication and universal
instance at $Sz$ under the base assumption, and discharge DI-case.
The closure composite is the unused-witness instance. External
structural induction on formulas, using the derived ordinary rules,
proves both formula composites. All these derivations are internal to
the respective theories; standard-model agreement alone is not the
argument.

The falsity bridge preserves contradictions in both directions.
The effective translations of finite derivations therefore establish
equiconsistency. They do not assert an internal proof of either
theory's own consistency.
\end{proof}

The proof is a written mathematical reduction with its inherited
arithmetic details reproduced below. The executable examples verify
particular proof certificates; they are not a Lean or Rocq verification
of this metatheorem.

% End included source


% Begin included source: theory/composition.tex
\section{A composed arity and predecessor checker}
\label{sec:composed-checker}
\label{sec:composition}

The construction in this section is deliberately a composition of two
existing mechanisms in the inspected drafts.  It changes the complete
certificate checker, not the underlying arithmetic foundation.  The main
program uses the strict list test \texttt{isinstance(a,list)} and has
612 normalized bytes, 266 significant Python tokens, and 310 AST nodes.
It is smaller than the current complete-program byte and token minima,
616 and 268, respectively.  It does not improve the separate 300-node
AST minimum.  Its actual source has 612 bytes and its ordinary gzip
level-9 encoding has 360 bytes, including the wrapper.  These figures
use CPython 3.12.14 and zlib 1.3.2.

The inherited arity checker~\cite{arity} supplies the sorted tree representation and
the logical branches.  The uniform-predecessor checker~\cite{uniform} supplies the
arithmetic packet interface.  Both have ordinary exit-status rejection
in this composition.  The underlying native calculus, numerical
bootstrap, reconstruction of ordinary intuitionistic logic, and
whole-domain bi-interpretation are the attributed inherited results.
Nothing here establishes an absolute minimum or a significant new
arithmetic foundation.

\subsection{The exact decoded language}

Numerical terms are variables, zero, successor, and initialized scalar
iteration
\[
 I(0,a;v.h)=a,\qquad
 I(Sn,a;v.h)=h[I(n,a;v.h)/v].
\]
The native formulas are equations and the paired binder; the checker
also uses conditional closures:
\[
 Q_y B\equiv B(0)\to\forall y\,B(y),\qquad
 D_{x,y}B\equiv\exists x\,Q_yB(x,y).
\]
The displayed equivalences give their ordinary realizations.  Their
primitive native rules are the inherited case/progress introductions,
closure use, paired elimination, equality transport, and computations.
Progress remains restricted to literal equations.

Canonical formulas are encoded by
\[
 s=t\longmapsto[s,t],\qquad
 D_{x,y}B\longmapsto[B],\qquad
 Q_yB\longmapsto[B,B,B].
\]
Only the third field of an arbitrary ternary formula has mathematical
content.  The repeated fields in a canonical closure are deliberate
redundancy.  Multiplication of a unary major array by three produces its
closure; multiplication of a binary or ternary major array by three
produces a high-arity formula, which denotes the fixed truth $0=0$.

The decoder is sorted.  A numerical root has protected-coordinate depth
three, whereas a formula root has depth zero.  Under $r$ protected
numerical slots the corresponding depths are $3+2r$ and $2r$.
At numerical root, coordinates $0,1,2$ are zero aliases; coordinates
$3+2j,4+2j$ are the two aliases of protected slot $j$.  Negative
coordinates admitted by the current floor denote rigid parameters.
Canonical inputs use integer coordinates.

A unary numerical node is successor.  Its child receives two dummy
zero slots.  A ternary numerical node $[n,a,h]$ is iteration.  Its
counter and initializer receive one dummy zero slot; its update
receives the actual state slot.  Empty numerical arrays, binary
numerical annotations, and numerical arrays of arity at least four
all denote zero.  At formula positions empty arrays, arrays of arity
at least four, and admitted scalar padding denote $0=0$.
These are fixed aliases, not extra predicate or function symbols.

The opener increments depth by
\[
                       5-\operatorname{len}(a)
\]
for every array.  The meaningful increments are $4$ for paired
binding or successor, $3$ for equality, and $2$ for closure binding
or iteration.  This is exactly the inherited arity codec.
At a meaningful scalar position it preserves an admitted coordinate
when $k\leq a<d$; otherwise, at odd depth it inserts its image after
checking that image at numerical depth three.  At even depth scalar
padding becomes truth.  A missing image at a used numerical hole
rejects.

\paragraph{The precise opening invariant.}
At a meaningful sorted root, a successful opener call implements
capture-free filling of the distinguished numerical hole under the
decoder just given.  The proof is structural induction on the raw tree.
For arities one, two, and three, the depth increment is exactly the
number of new protected slots plus any required sort change.  At
ignored annotations and surplus-arity nodes the interpretation is
constant, so traversal below them cannot introduce a meaningful free
variable.  The inserted image is checked at its own numerical root;
its free parameters are rigid negative labels, so inserting it below
protected slots captures none of them.

This is a conditional semantic invariant.  It is not a claim that
opening is idempotent on every raw tree, nor that every semantic alias
must accept the same certificate.  For example, opening a high-arity
truth array can alter its ignored children; a later traversal may
reject those children without contradicting its constant decoded
meaning.  Canonical object expressions avoid surplus nodes, while
administrative high-arity truth exports are justified separately.

\subsection{An exact reduction of the arithmetic branch}

A new arithmetic instruction is
\[
                         [\mathtt{null},p,a,h].
\]
If the raw predecessor $p$ is truthy, the checker constructs
\[
 H=\bigl[\,[[p],a,h],\ i(h,[p,a,h],3)\,\bigr].
 \tag{C1}\label{eq:composed-step}
\]
If $p$ is falsy, it constructs
\[
                         H=[[p,a,h],a].
 \tag{C0}\label{eq:composed-base}
\]
It then allocates its next parameter and opens this local hypothesis
with image zero.  The test selecting the arithmetic branch is the
literal null test.  Literal zero therefore remains available as a
paired-introduction witness.

There is a particularly direct soundness reduction, which also checks
the change of numerical opening depth.  The old arity checker accepts
constructed-computation packets of the form
$[m,s,u,t]$, where $s=[n,b,r]$ and falsy $t$ selects computation.
Use the following old packets for the new packet above:
\[
\begin{array}{c|c}
\text{Condition}&\text{Corresponding old arity packet}\\\hline
p\text{ truthy}&[\,[],\,[p,a,h],\,1,\,0\,]\\
p\text{ falsy}&[\,[],\,[0,a,h],\,p,\,0\,].
\end{array}
\tag{Packet correspondence}\label{eq:packet-correspondence}
\]
The empty-array marker is essential in this comparison.  In the common
continuation, both new marker null and old marker $[]$ make
\texttt{len(goal) \% 2 == marker} false.  Both therefore select image
zero for the supplied hypothesis, independently of the goal's arity.
One must not compare with the old numeric-zero marker and silently
assume its image choice is unchanged.

For truthy $p$, the old destructuring is
$n=p,b=a,r=h,u=1$.  Its constructed source and recursive-image opening
are exactly \eqref{eq:composed-step}, including the opening at depth
three.  For falsy $p$, Python's expression $u\mathbin{\mathrm{and}}[n]$
returns $u=p$, so the old source is exactly
\eqref{eq:composed-base}, retaining even the same raw zero alias.
The unused first field of the old triple may be zero.  Thus the generated
raw equation, old-floor image checks, next floor, opened goal, opened
hypothesis, and continuation context are identical in the two
executions.  No appeal to a different meaning of a raw numerical
coordinate is needed.

Consequently the inherited arithmetic-supply proof applies without
alteration.  In the state-dependent successor case, the call
$i(h,[p,a,h],3)$ checks its recursive image at numerical root depth
three whenever the state is used in a relevant position.  This checks
the predecessor, initializer, and update at the old floor.  Coordinates
three and four are the state aliases.  An escaping relevant coordinate
cannot be hidden by deferring that image check.  If the update does not
use the state, the successor equation has its ordinary constant-update
form and does not require the recursive image.  The source itself is
still traversed when the supplied equation is opened.

In the base case, every successfully admitted falsy numerical counter
in the strict-list profile denotes zero: the relevant aliases are zero,
false, numerically equal zero values, and the empty array.  Null,
strings, and dictionaries reject when used in a meaningful numerical
position.  This gives an ordinary zero-computation instance.

\paragraph{Fresh parameters and the zero-instance argument.}
The exact reduction above is useful because it prevents a stronger
arithmetic axiom from entering through the proof interface.  The
inherited supply argument permits an equation $H(\vec z)$ when its
simultaneous zero instance is a computation theorem and $\vec z$ is
fresh for the retained context and goal.  If
\[
 \Gamma,H(\vec z)\vdash F\quad\hbox{and}\quad\vdash H(\vec0),
\]
substituting zero throughout the first derivation leaves $\Gamma,F$
unchanged; cutting in the second derivation gives $\Gamma\vdash F$.
This is a finite native proof transformation, not a consistency or
reflection assumption.

When the floor drops from $K$ to $K-1$, every relevant new label in
$[K-1,K)$ is covered by this simultaneous substitution.  For canonical
integer syntax only $K-1$ is new.  If fractional JSON numbers are
accepted as raw aliases, they are additional distinct fresh labels;
they must not silently be identified with the designated integer
parameter.  The known-zero-instance argument covers the finitely many
that occur.  Previously checked goals and retained hypotheses contain
none of these names in meaningful positions.  The composition does not
need to assert that every raw supplied equation is uniformly a theorem
in every newly admitted label.

\subsection{The logical branches and zero witnesses}

Every truthy-major branch is inherited unchanged from the strict arity
checker.  The major is proved first.  An exact equation endpoint permits
transport, while an exact triplicate closure permits specialization
only when the old endpoint is numerical zero.  The explicit
\texttt{not source} condition is retained.  The continuation either
exports the closure of a paired major at a fresh parameter, or exports
a formula with the fixed truth interpretation.

Every remaining falsy non-null marker selects paired introduction.
The expression \texttt{i(*goal, witness, 2)} enforces the unary shape of
a canonical paired formula.  Binary equation goals and ternary closure
goals supply too many arguments and reject.  An empty goal has the
constant-truth interpretation, so any accepted administrative empty-goal
case is harmless.  Surplus-arity goals also fail this call.

For a canonical goal $D_{x,y}B$, opening at depth two preserves $y$ and
fills its witness coordinate with the supplied term $t$.  The checker
then proves the successor instance at a fresh $z$.  Marker $[]$
provides $B(t,0)$ as the local assumption.  Marker $0$ provides
$B(t,z)$ only when the opened matrix is a binary, literal equation;
for a compound matrix it provides the zero instance instead.  These
are precisely the inherited case and equation-progress rules.

This argument allows the actual numerical witness zero.  It does not
identify zero with a missing witness.  A missing image is rejected when
the witness coordinate is semantically used.  If the coordinate is
unused, zero is a legitimate mathematical representative.  The
canonical completeness argument needs only ordinary arrays and integer
coordinates, so the additional harmless raw aliases do not increase
its logical requirements.

\subsection{Completeness, including the zero predecessor}

Canonical encoding of the native logical rules follows the inherited
arity proof.  The only change needed in arithmetic is to handle the
syntactic falsity of a literal-zero predecessor in a successor
computation.
Define the finite term
\[
                         B(p)=I(0,p;w.0).
\]
The base computation proves $B(p)=p$, and the syntax of $B(p)$ is a
nonempty array.  For any update $h$ with state variable $v$, put
\[
 F(z)\;:\equiv\;
 I(Sz,a;v.h)=h[I(z,a;v.h)/v].
\]
The new arithmetic branch directly proves $F(B(p))$ because the
predecessor $B(p)$ has truthy syntax.  A second arithmetic instruction
proves $B(p)=p$.  Equality transport through the full formula template
$F(z)$ then proves $F(p)$.  This is an explicit three-node derivation:
one boxed successor computation, one zero computation, and one
transport.  It works in particular for $p=0$, and works for every
well-scoped predecessor if a uniform compiler is preferred.

The zero computation itself is direct: use predecessor zero.  All
other truthy canonical predecessors may use direct successor packets.
The inherited eigenparameter renaming and continuation serialization
then compile every finite native proof.  The supplied untrusted
producer implements this box expansion, combines the exact arity
numerical/formula codecs with the uniform-predecessor emitter, and
retains the named-context discipline.  A checker never imports that
producer or any arithmetic evaluator.

It follows that the canonical consequence relation is exactly the
inherited native calculus.  Its ordinary realization is therefore
whole-domain bi-interpretable and equiconsistent with iPRA by the
already supplied logical and arithmetic bridge.  The interpretations
and their identity numerical round trips are unchanged.  This transfer
is a consequence of equality of the presented calculus, not an inference
from finite testing or from agreement on standard numerals.

\subsection{Ancillary profiles and the exact validation boundary}

The byte profile composes the inherited exact-JSON list predicate and
short-circuit endpoint expression.  On canonical object formulas the
replacement endpoint is a nonempty array, so this expression has the
same behavior as the explicit conditional.  On a raw falsy truth alias
it can discard the conversion and retain the old goal.  It then requires
a proof of that old goal and does not introduce an unsound implication.
The profile measures 597 bytes, 267 tokens, and 315 AST nodes; its
actual gzip-9 stream has 356 bytes.

The AST profile uses the inherited exception-based opener and fused
initialization.  It measures 639 normalized bytes, 267 tokens, and
298 AST nodes, with a 360-byte gzip-9 stream.  Nonempty malformed strings
reject through Python's finite recursion limit; an ideal unbounded
interpreter could diverge on those malformed inputs.  This qualification
is retained, and the profile is not the recommended total-input parser.
An immutable-context continuation loop gives 296 AST nodes at the cost
of 689 normalized bytes.

All measurements include input parsing calls, scope traversal, logical
and arithmetic verification, and the acceptance driver.  The shared
Python interpreter and JSON library remain outside the source counts,
as in the inspected drafts.  The measured primary program imports only
\texttt{json}.  Ordinary nonzero process exit rejects; acceptance requires
exit zero and the exact line \texttt{ACCEPT}.

The delivered finite suite regenerates the canonical proof examples,
compares the producer against its documented frozen-source adaptation,
checks the entire public programs normally and with Python optimization,
and checks malformed-input and guard-removal controls.  The false
statement controls include $0=1$, $I(1,1;v.v)=0$, and
$I(0,1;v.0)=0$.  Separate controls test a forbidden compound-induction
step and a free parameter in a requested closed goal.  These controls
are evidence about the implementation; the universal correctness and
bi-interpretability claims rest on the written arguments and inherited
bridge.  No Lean or Rocq formalization is claimed.

\subsection{The entire selected executable}
\begin{lstlisting}[language=Python]
import json
k = 0
i = lambda a, t=None, d=0: [i(b, t, d + 5 - len(a)) for b in a] if isinstance(a, list) else a if k <= a < d else d % 2 and i(t, d=3)

def c(a, g):
    global k
    if i(a) in g:
        return
    m, s, u, t = q()
    if m:
        c(m, g)
        if m in ([s, u], not s and [t] * 3) and i(t, u) == a:
            a = i(t, s)
        h = m * 3
    elif m is None:
        h = [[s and [s], u, t], i(t, [s, u, t], 3) if s else u]
    else:
        h = a = i(*a, t, 2)
    k -= 1
    c(i(a, [k]), g + [i(h, k if m or len(a) % 2 == m else 0)])
q = json.load(open(0)).pop
c(q(), [])
print('ACCEPT')
\end{lstlisting}

% End included source


% Begin included source: theory/conjunctive_obstruction.tex
% New contribution relative to the inspected 2026-11-10 draft corpus.
% This fragment is intended for inclusion after standard theorem environments.
\section{Conjunctive induction still leaves a symmetry obstruction}
\label{sec:pp-obstruction}

The earlier eventually-constant-power obstruction in the inspected
\texttt{ipra\_research\_2026-11-10} draft, specifically
\texttt{theory/finite\_parameter\_obstruction.tex}, includes equation
induction, finite conjunctions of equations, and universal Horn
arithmetic laws~\cite{research}.  That source already attributes the finite-parameter
involution argument to earlier unary-amalgam work; neither the power
construction nor the automorphism method is claimed here as new.
The result below first extends that particular construction to existential
witnesses: both the induction matrices and the regular axioms may contain
positive existential quantification.  A further closure theorem allows
unrestricted finite alternation of existential and universal quantifiers,
provided that the formulas use conjunction and no disjunction,
implication, or negation internally.  Both layers extend the existing
construction; neither is presented as a literature-wide priority claim
or a positive axiomatization result.

\subsection{Exact hypotheses}

Let $\mathcal N$ be any one-sorted finitary expansion of the standard
structure $(\mathbb N,0,S)$ by total functions and relations.  The
signature can be infinite; each symbol has finite arity, and each formula
is finite.  In particular, it can contain every primitive-recursive
function.  A \emph{primitive-positive} formula is built from atomic
formulas and $\top$ by finite conjunction and existential quantification.
Neither disjunction, implication, negation, nor universal quantification
is permitted inside such a formula.

We permit two kinds of regular axioms:
\[
 \forall\vec x\,\bigl(\varphi(\vec x)\to\psi(\vec x)\bigr),
 \qquad
 \forall\vec x\,\bigl(\varphi(\vec x)\to\bot\bigr),
 \tag{Reg}\label{eq:pp-regular}
\]
where $\varphi,\psi$ are primitive-positive and the displayed sentence
is true in $\mathcal N$.  The second form explicitly includes negative
constraints such as $0=S0\to\bot$.  Terminology differs on whether
regular logic includes falsity, so this form is listed separately.

For every primitive-positive $A(n,\vec p)$, include the full induction
sentence
\[
 \forall\vec p\left[
 \left(A(0,\vec p)\wedge
       \forall n\,(A(n,\vec p)\to A(Sn,\vec p))\right)
 \to \forall n\,A(n,\vec p)\right].
 \tag{PPI}\label{eq:pp-induction}
\]
Parameters $\vec p$ remain fixed during induction.  This is equivalent
to the usual contextual induction rule for this class of matrices in
intuitionistic first-order logic.  Retained assumptions can contain
arbitrary formulas.  The claim below does not restrict them to
primitive-positive formulas.

Let $T_{\mathrm{pp}}(\mathcal N)$ be intuitionistic first-order logic
with equality, all sentences \eqref{eq:pp-regular}, and all instances
\eqref{eq:pp-induction}.  This semantic upper theory need not have a
recursive axiom set.  The obstruction also applies to any recursively
axiomatized subtheory containing the same symbols and using only axioms
from these classes.  More generally, it applies to any theory satisfied
by the model constructed below.

\subsection{The model and its positive existential formulas}

Write
\[
 E=\{a\colon\mathbb N\to\mathbb N:
       a\text{ is eventually constant}\}.
\]
Interpret every function of $\mathcal N$ coordinatewise.  For an
$r$-ary relation $R$, put
\[
 E\models R(a_1,\ldots,a_r)
 \quad\Longleftrightarrow\quad
 \mathcal N\models R(a_1(i),\ldots,a_r(i))
       \text{ for every }i\in\mathbb N.
\]
Equality is actual equality of sequences, and constants are constant
sequences.  Finitarity ensures that coordinatewise functions take
eventually constant inputs to eventually constant outputs.  Denote the
resulting structure by $\mathcal E$.

\begin{lemma}[Primitive-positive transfer]
\label{lem:pp-transfer}
For every primitive-positive $\varphi(\vec x)$ and every finite tuple
$\vec a$ from $E$,
\[
 \mathcal E\models\varphi(\vec a)
 \quad\Longleftrightarrow\quad
 \mathcal N\models\varphi(\vec a(i))
                 \text{ for every }i\in\mathbb N.
 \tag{PP}\label{eq:pp-transfer}
\]
\end{lemma}
\begin{proof}
Use external induction on the fixed finite formula.  Atomic formulas
and finite conjunctions follow immediately from the coordinatewise
definitions.  The forward existential direction follows by evaluating
a witness at each coordinate.

For the reverse existential direction, suppose
$\mathcal N\models\exists y\,\theta(y,\vec a(i))$ at every
coordinate, and assume the transfer statement for $\theta$.  There is
a finite set $F$ such that $\vec a(i)=\vec a_\infty$ for every
$i\notin F$.  Choose a witness $b_\infty$ for the common tail tuple,
and a witness $b_i$ separately for each of the finitely many coordinates
$i\in F$.  The sequence with values $b_i$ on $F$ and $b_\infty$
elsewhere belongs to $E$.  The induction hypothesis gives
$\mathcal E\models\theta(b,\vec a)$, hence the required
existential statement.  Only finitely many witness choices are used.
Nested or tuple existential quantifiers are handled by repeated use of
this argument.
\end{proof}

\begin{lemma}[Regular axioms and primitive-positive induction]
\label{lem:pp-model}
The classical structure $\mathcal E$ satisfies every axiom of
$T_{\mathrm{pp}}(\mathcal N)$.
\end{lemma}
\begin{proof}
For a regular axiom, a true antecedent in $\mathcal E$ gives a true
antecedent at every coordinate by Lemma~\ref{lem:pp-transfer}.  Validity
in $\mathcal N$ gives the conclusion at every coordinate; the reverse
direction of the lemma reconstructs its witnesses in $\mathcal E$.
For a negative regular constraint, a true antecedent in $\mathcal E$
would give a true antecedent in $\mathcal N$ at coordinate zero,
contrary to the constraint's validity there.

Fix an instance of \eqref{eq:pp-induction} and a parameter tuple
$\vec p\in E$.  Suppose its base and progress hypotheses hold in
$\mathcal E$.  Fix a coordinate $i$.  The transfer lemma gives
$\mathcal N\models A(0,\vec p(i))$.  If
$\mathcal N\models A(k,\vec p(i))$, form the sequence $a$
having value $k$ at coordinate $i$ and value $0$ elsewhere.  At
coordinate $i$, $A(a(i),\vec p(i))$ holds by the supposition;
at every other coordinate it holds by the transferred base hypothesis.
The reverse transfer lemma therefore gives
$\mathcal E\models A(a,\vec p)$.  Apply the assumed progress
sentence in $\mathcal E$, and then transfer back at coordinate $i$.
The result is
$\mathcal N\models A(k+1,\vec p(i))$.

Ordinary external induction on the standard number $k$ now gives
$\mathcal N\models A(k,\vec p(i))$ for every $k$.  This holds
at each coordinate $i$.  For any $a\in E$, the reverse transfer
lemma gives $\mathcal E\models A(a,\vec p)$.  Hence the whole
induction instance holds.  The parameter tuple has remained fixed
throughout.  Nothing here assumes that witnesses for $A$ are terms or
are uniformly definable functions of $n$.
\end{proof}

A classical model is a one-node intuitionistic Kripke model.  Thus the
lemma supplies a model of the intuitionistic theory, not merely evidence
about its classical restriction.  In particular, any derivation using
arbitrarily quantified retained assumptions must be valid in this model.

For orientation, successor coverage already fails on the original
domain.  The sequence equal to $1$ at one coordinate and $0$ elsewhere
is not the zero sequence and is not a successor, since every successor
sequence has a positive value at every coordinate.  This observation
alone would not exclude an interpretation on a different definable
domain; the following symmetry argument is what handles that freedom.

\subsection{The automorphism obstruction}

\begin{lemma}[Finite parameters leave an involution]
\label{lem:pp-involution}
For every finite tuple $\vec p\in E$, the structure
$(\mathcal E,\vec p)$ has a nonidentity automorphism of order two.
\end{lemma}
\begin{proof}
Outside a finite set, every component of $\vec p$ is constant.  Choose
two distinct coordinates $i,j$ outside this set and swap them.
Coordinatewise interpretation makes the swap an automorphism preserving
all function and relation symbols, irrespective of the size of the
signature.  It fixes $\vec p$ and has order two.  It is nonidentity
because the sequence equal to $1$ at $i$ and $0$ elsewhere is moved to
the corresponding sequence supported at $j$.
\end{proof}

\begin{theorem}[No ordinary bi-interpretation with iPRA]
\label{thm:pp-obstruction}
$T_{\mathrm{pp}}(\mathcal N)$ is not bi-interpretable with iPRA by
ordinary finitary first-order syntactic interpretations, of finite
dimension, whose composites are provably definably isomorphic to the
respective identity interpretations.  Arbitrary formula-defined domains,
relation/function graphs, and quotient equality are allowed.  Each
model may use a fixed finite parameter tuple, including parameters for
the return isomorphisms, as long as the claimed interpretations and
isomorphisms are required to exist in that model.  In particular, this
covers a syntactic parameter scheme whose finite parameter condition is
provably inhabited and under which all these facts are proved uniformly.
The same conclusion holds for every theory having $\mathcal E$ as a
model.  A parameter condition that simply excludes $\mathcal E$ and is
not proved inhabited there is outside the claim.
\end{theorem}
\begin{proof}
Every classical model of iPRA has a definable linear order.  For example,
in the standard primitive-recursive language use
$x\le y\leftrightarrow\exists z\,(x+z=y)$; iPRA proves that this
is a linear order.  An order-preserving automorphism of finite order is
the identity: if $a<\sigma(a)$, iterating the strict inequality gives
$a<\sigma(a)<\cdots<\sigma^r(a)=a$, and the case
$\sigma(a)<a$ is analogous.

Suppose there were a bi-interpretation.  Apply the interpretation to
$\mathcal E$ to obtain a classical iPRA model $\mathcal M$.
Collect the finitely many parameters of both interpretations and their
definable return isomorphisms; representatives of any parameters in an
interpreted quotient add only finitely many elements of $E$.
Lemma~\ref{lem:pp-involution} supplies a nonidentity involution
$\sigma$ fixing all of them.

Every first-order interpretation carries such an automorphism to the
interpreted structure, including when its domain is a definable set of
tuples modulo a definable equivalence relation.  The induced
automorphism $\widehat\sigma$ of $\mathcal M$ has order dividing
two, so the preceding linear-order argument makes it the identity.
Its further interpretation is therefore the identity on the
double-interpreted structure.  The source-side return isomorphism is
definable with parameters fixed by $\sigma$, hence is equivariant:
\[
 \eta(\sigma(a))=\widehat{\widehat\sigma}(\eta(a))=\eta(a).
\]
Since $\eta$ is injective, $\sigma(a)=a$ for every $a\in E$,
contrary to the choice of $\sigma$.
\end{proof}

\begin{corollary}[Additional sorts do not evade this obstruction]
\label{cor:pp-many-sorted}
Let $\mathcal A$ be a finitary many-sorted structure with nonempty
sorts and a distinguished numerical sort that is the standard
$(\mathbb N,0,S)$.  Other sorts may, for example, consist of finitary
objects or specified numerical functions, with arbitrary total finitary
operations and relations between the sorts.  Form the intuitionistic
theory of all true many-sorted regular sequents and negative regular
constraints, together with induction on the numerical sort for every
primitive-positive matrix.  It is not bi-interpretable with iPRA by
the ordinary syntactic interpretations of
Theorem~\ref{thm:pp-obstruction}.  Interpretation domains may use finite
products of any of the source sorts and definable quotients.
\end{corollary}
\begin{proof}
For each sort with domain $A_s$, take the eventually constant sequences
$\mathbb N\to A_s$, and interpret all symbols coordinatewise, using
the same coordinate index set for every sort.  A fixed formula and its
parameter tuple involve only finitely many sorts and coordinates with
exceptional values.  The primitive-positive transfer proof therefore
works unchanged: existential witnesses have their specified sorts,
and finite exceptional witnesses plus one common tail witness give
an eventually constant witness of that sort.  This proves every true
regular axiom.

For induction, the varying sequence $a_{i,k}$ belongs to the
distinguished numerical sort.  All other parameter sorts stay fixed.
The proof of Lemma~\ref{lem:pp-model} consequently proves every
primitive-positive numerical-induction instance, even when witnesses
and parameters range over other sorts.

Swap two coordinates outside the combined finite support of all
interpretation and return-map parameters, simultaneously in every
sort.  This is an automorphism of the whole many-sorted structure.
It is nonidentity already on the numerical sort, by the sequence
supported at one of the swapped coordinates.  Its induced
automorphism on an interpreted iPRA model must be the identity.
Definability of the source return isomorphism, on every source sort,
then makes the original automorphism the identity on every sort.
The numerical-sort contradiction suffices.  No faithfulness of the
action on each other sort is assumed or needed.
\end{proof}

\subsection{Universal quantifiers give a stronger closure theorem}

The primitive-positive result is not the limit of this method.  Let
$\mathcal C$ be the class of \emph{conjunctive formulas} generated from
atomic formulas and $\top$ by $\wedge$, $\exists$, and $\forall$.
Quantifier alternation is unrestricted.  There is still no disjunction,
implication, or negation inside a matrix from $\mathcal C$.
The explicit grammar is the convention used here; no identification
with a differently defined syntactic fragment is assumed.

\begin{theorem}[Conjunctive transfer and induction]
\label{thm:conjunctive-induction}
For every $\varphi(\vec x)\in\mathcal C$,
\[
 \mathcal E\models\varphi(\vec a)
 \quad\Longleftrightarrow\quad
 \mathcal N\models\varphi(\vec a(i))
                  \text{ for every }i\in\mathbb N.
\]
Consequently $\mathcal E$ satisfies induction for every matrix in
$\mathcal C$, as well as every sentence true in $\mathcal N$ of the
forms
\[
 \forall\vec x\,\bigl(\varphi(\vec x)\to\psi(\vec x)\bigr),
 \qquad
 \forall\vec x\,\bigl(\varphi(\vec x)\to\bot\bigr),
 \qquad \varphi,\psi\in\mathcal C.
\]
These conclusions hold also in the many-sorted setting of
Corollary~\ref{cor:pp-many-sorted}, with induction in the numerical
sort.  The resulting stronger theories are still not bi-interpretable
with iPRA under the interpretations specified in
Theorem~\ref{thm:pp-obstruction}.
\end{theorem}
\begin{proof}
Induct on the fixed finite formula.  The atomic and conjunction cases
are unchanged.  The existential proof of
Lemma~\ref{lem:pp-transfer} uses only that the body already satisfies
transfer and that its finite parameter tuple has a common tail.
It therefore works even when the body contains universal quantifiers.

For the universal case, suppose
$\mathcal E\models\forall y\,\theta(y,\vec a)$.  Fix a
coordinate $i$ and any element $b$ of the corresponding standard sort.
The constant sequence $\overline b$ is an element of the sequence
sort.  Instantiate the universal statement at $\overline b$ and apply
the induction hypothesis at coordinate $i$.  This yields
$\mathcal N\models\theta(b,\vec a(i))$.  As $b$ was arbitrary,
the coordinate satisfies the universal statement.  Conversely, if
every coordinate satisfies that universal statement, then for every
sequence $b$ every coordinate satisfies $\theta(b(i),\vec a(i))$.
The induction hypothesis gives
$\mathcal E\models\theta(b,\vec a)$ for every such $b$,
which is the required universal statement in $\mathcal E$.

The proofs of regular-axiom preservation and induction in
Lemma~\ref{lem:pp-model} use only exact coordinatewise transfer;
apply them to this larger class.  The finite-parameter involution and
return-isomorphism obstruction use only the unchanged underlying model,
so they apply verbatim.  The proof allows arbitrarily alternating
quantifiers; it supplies no term-witness or decidability assumption.
\end{proof}

The one-sorted primitive-positive and many-sorted statements above are
thus special cases.  Even unbounded quantifier alternation in this
conjunctive fragment does not force the uniform branch information
missing from successor coverage.

\paragraph{Function objects and extensionality.}
The many-sorted statement also applies when a sort consists of specified
actual numerical functions and the signature includes a total finitary
application operation. Extensionality has the allowed form
\[
 \bigl(\forall n\;\operatorname{app}(f,n)=\operatorname{app}(g,n)\bigr)
       \to f=g.
\]
Its antecedent is conjunctive, so the sequence model preserves it.
Thus adding function objects, application, and extensionality does not
by itself evade the obstruction. Equational combinator laws and any
other standard-true conjunctive sequents may be included as well.
This observation does not address a calculus with additional principles
outside the stated class; it identifies the exact assumptions under
which that possible change of domain still retains the symmetry.

\paragraph{Scope.}
The theorem is about ordinary first-order bi-interpretations, not
equiconsistency alone.  It allows existentially specified witnesses and
conjunctive induction of arbitrary quantifier alternation, so changing
a term-defined operation to a positive existential graph does not by
itself evade this obstruction.
It does not rule out arbitrary one-way interpretations, nonregular
logical rules, disjunctive coverage axioms, or arbitrary formula-defined
expansions that are not interpreted pointwise.  Nor does it yield a
numerical lower bound for proof-checker source size.  Its practical use
is to rule out an entire family of proposed positive existential
completions before attempting a full bi-interpretation proof.

% End included source


% Begin included source: theory/validation.tex
\section{Executable evidence, attribution, and research limits}
\label{sec:validation}

\subsection{What the delivered checks establish}

The selected kernel and its four explanatory or ancillary counterparts
were checked on 31 named proof and adversarial fixtures, under ordinary
Python and Python's \texttt{-O} optimization mode. The resulting 310
complete-program executions matched their expected outcomes. Forty
malformed-input checks and 24 targeted guard-control executions also
passed. Successful checking requires both process exit zero and exactly
the line \texttt{ACCEPT}; a traceback alone is not a verdict convention.

The positive examples include zero computation with an arbitrary seed,
successor computation using the previous state, retained parameters,
transport, closure specialization, paired introductions and eliminations,
zero witnesses, and contextual induction. The fixture producer is
untrusted and is not imported by a checker. It is independently
regenerated from two frozen upstream producers, and the generated
packets are compared with the delivered fixtures before checking.

The negative controls target actual failure modes. Removing selected
guards makes deliberately changed kernels accept false equations such
as $0=1$, $I(1,1;v.v)=0$, or $I(0,1;v.0)=0$. Other controls exercise
an attempted induction on a compound formula and an unallocated
parameter in a requested closed goal. These are targeted checks of the
trust boundary. They are not a proof that no other implementation defect
exists, and they are not a formal verification of the metatheorems.

The ordinary-list primary traversal handles finite JSON trees by the
specified checking operations, subject to physical resource bounds.
The ancillary exception-based opener has the additional qualification
that rejection of certain malformed nonempty strings uses the runtime's
recursion limit. Large valid proofs can exceed ordinary memory or stack
limits in every delivered implementation. The mathematical completeness
statement is about unbounded finite derivations, not an assertion that a
fixed physical process accepts every arbitrarily large proof.

\subsection{Exact examples for the symmetry theorem}

The file \file{code/eventually_constant.py} represents an eventually
constant sequence by its constant tail and a finite exception map.
Its output illustrates the infinite model exactly; it does not replace
the transfer and induction proofs by bounded experiments.
Let $e_i$ be one at coordinate $i$ and zero elsewhere. Then:
\begin{center}
\begin{tabular}{lll}
\toprule
Property & In the sequence model & Reason\\
\midrule
$e_i=0$ & False & Coordinate $i$ is one\\
$\exists y\;(e_i=Sy)$ & False & All other coordinates are zero\\
$\operatorname{sg}(e_i)=0$ & False & Its $i$th coordinate is one\\
$\operatorname{sg}(e_i)=1$ & False & Its tail is zero\\
$e_i\le e_j$, $i\ne j$ & False & Fails at coordinate $i$\\
$e_j\le e_i$, $i\ne j$ & False & Fails at coordinate $j$\\
\bottomrule
\end{tabular}
\end{center}

Here $\operatorname{sg}(0)=0$ and $\operatorname{sg}(Sn)=1$,
interpreted coordinatewise. All its true conjunctive identities still
hold. The failure is uniform disjunctive information: being Boolean
at every coordinate does not mean being globally the constant zero
or the constant one. This explains why standard Boolean coding cannot
silently turn the model's conjunctive induction into full quantifier-free
induction with the usual logical equivalences.

The example also swaps two coordinates outside the finite supports of
specified parameters, fixes those parameters, and moves $e_i$ to $e_j$.
The theorem proves this for every finite tuple and every finitary
pointwise signature. The program merely makes representative instances
easy to inspect.

\subsection{What was investigated but not claimed}

Two broader positive directions were examined. First, replacing
computation supply with numerical definitional conversion would remove
an explicit branch, but a correct reducer must respect the codec's
dummy slots, state binder, neutral counters, and rigid parameters. Its
extra machinery did not yield a smaller implemented checker. Moreover,
a uniformly total primitive-recursive evaluator for arbitrary
primitive-recursive descriptions is impossible by the usual diagonal
argument. A recursive normalizer or a bounded conversion certificate
is a different option and is not excluded by that observation.

Second, genuine binders in the proof-certificate grammar could let the
opener allocate unforgeable eigenparameters, replacing predictable
global negative labels. The flat instruction stream currently needs
those labels so that later instructions can refer to earlier
eigenparameters. A new scoped grammar would therefore require a
corresponding proof-opening discipline. No complete smaller construction
was obtained along that route.

Deleting DI-case for literal equations is not justified. An ordinary
conditional closure can hold because its zero antecedent is false,
without its matrix being progressive at every argument. Reinterpreting
atomic paired formulas as progressiveness also requires a matching
elimination rule and a proof of the source-side formula composite.
These obligations cannot be discharged just by recovering some ordinary
logical macros. The existing drafts already analyze this difficulty;
it is not claimed as a new negative result here.

\subsection{Reproduction and the remaining target}

From a fresh extraction, with Python 3.12 and a normal PDFLaTeX
installation:
\begin{lstlisting}
python3 reproduce.py
python3 code/checker.py < tests/fixtures/I0_arbitrary_seed.json
python3 code/checker.py < tests/fixtures/contextual_induction_seed_congruence.json
python3 build.py
\end{lstlisting}
The first command verifies payload hashes, recomputes the complete-module
and gzip measurements, checks producer provenance, runs the delivered
proof fixtures and guard controls, and reproduces the exact sequence
examples. It needs no network. The last command rebuilds the report
without shell escape. A standalone LaTeX source is supplied as well as
the modular sources.

The package preserves frozen comparison programs and the exhaustive
source inventory. Repeating the complete corpus scan additionally needs
the thirteen original archives from the pinned GitHub revision;
\file{reference/census.py} records the scan procedure. The regular
reproduction command uses the included frozen programs and requires
none of the original ZIP downloads.

The outcome is intentionally narrower than an absolute solution to the
research problem. There is a complete implementation below the supplied
byte and token envelope, and a newly extended obstruction that rules
out an otherwise tempting family of changed-language and changed-domain
approaches. A significant new positive axiom basis, a checker below all
three structural minima in one program, and a global simplicity theorem
remain unproved.

% End included source


\appendix
\clearpage
\section*{Attribution and the complete inherited bridge}
\addcontentsline{toc}{section}{Attribution and the complete inherited bridge}
The next five sections are reproduced unchanged from the
\file{ipra_uniform_predecessor_2026-11-10.zip} archive at the pinned
commit~\cite{snapshot,uniform}. They contain the logical reconstruction,
the ordinary first-order comparison theory, the scalar predecessor,
actual equality reflection, and the primitive-recursion compiler.
They are included to make the asserted iPRA equivalence inspectable in
this package. These results are inherited and are not claimed as new here.
Their exact hashes and original member paths are recorded in
\file{evidence/inherited_provenance.json}. The repository's MIT No
Attribution license is included in \file{reference/PROVEIT_LICENSE.txt}.

The clause notation in the first appendix is the older presentation of
the same conditional closure used as $\Q$ in the main text. A clause
$\langle B(0);y.B(y)\rangle$ is represented by $\Q_y B$; an ordinary
assumption $F$ is a constant clause $\langle H;v.F\rangle$, with
$H$ a fixed computation theorem, and corresponds to the formula assumption
$F$, not to the tautological closure $\Q_vF$. Conversely $\Q_y B$ is interderivable
with $\D_{x,y}B$ for an unused $x$. The public-$\Q$ presentation in the
main text and the appendix's clause presentation therefore have the same
closed native theorems. This interface is also used by the frozen
upstream checkers; it is not an extra logical oracle in the new program.


% Begin included source: reference/inherited_choice_basis.tex
% New mathematical result: complementary introduction choices.
% Requires amsmath/amssymb/amsthm and theorem, lemma, proposition environments.
% The inherited arithmetic-to-iPRA bridge must be supplied separately.

\section{Complementary introductions with one local assumption}
\label{sec:choice-introduction}

The old paired introduction retains both the zero instance and an
induction instance when its matrix is an equation. Both together are
unnecessary. A certificate may choose either one, provided the induction
choice is restricted to a literal equation. This changes neither the
formula language nor the meaning of a formula. The omitted combination
is recovered by a short nested introduction--elimination construction.

\subsection{Complete native syntax and rule inventory}

The scalar numerical language is
\[
 t::=x\mid0\mid St\mid I(n,a;u.h).
\]
Only $h$ binds the state variable $u$; $n$ and $a$ do not.
The update may contain fixed numerical parameters. The computation
axioms in every retained context are
\[
 I(0,a;u.h)=a,\qquad
 I(Sn,a;u.h)=h[I(n,a;u.h)/u].
 \tag{I0, IS}\label{eq:choice-computations}
\]
The argument also applies unchanged to the full fixed-description PR
language with its projection, composition, and recursion equations.
No numerical variable ranges over function descriptions, and no uniform
evaluator is postulated.

Native formulas and reusable context clauses are
\[
 A::=(s=t)\mid\mathsf D_{x,y}A,\qquad
 \langle A;y.B(y)\rangle,\quad y\notin\operatorname{FV}(A).
\]
Bound variables are renamed before capture-free substitution.
The two variables $x,y$ in a paired formula are distinct.
An ordinary assumption $F$ is the constant clause
$\langle H;v.F\rangle$, where $v$ is absent from $F$ and $H$ is the
fixed computation theorem $I(0,0;u.0)=0$. Thus ordinary hypothesis
use is a clause Use preceded by this computation proof.

There are computation axioms and exactly the following inference
patterns:
\[
\frac{\Gamma\vdash A}{\Gamma\vdash B(t)}
\quad\mathsf{Use}
\quad(\langle A;y.B(y)\rangle\in\Gamma),
\]
\[
\frac{\Gamma\vdash\mathsf D_{x,y}B(x,y)\qquad
      \Gamma,\langle B(a,0);y.B(a,y)\rangle\vdash F}
     {\Gamma\vdash F}
\quad\mathsf{DE},
\qquad
\frac{\Gamma\vdash s=t\qquad\Gamma\vdash F(s)}
     {\Gamma\vdash F(t)}
\quad\mathsf{Tr},
\]
\[
\frac{\Gamma,B(t,\eta_\varepsilon(z))\vdash B(t,Sz)}
     {\Gamma\vdash\mathsf D_{x,y}B(x,y)}
\quad\mathsf{DI}_{\varepsilon},
\qquad
\eta_0(z)=0,\quad\eta_1(z)=z.
\tag{Choice DI}\label{eq:choice-di}
\]
Here $\varepsilon$ is a certificate choice in $\{0,1\}$, and
$\varepsilon=1$ is permitted only when $B$ has a literal equation
root. An equation beneath a binder is not a literal equation.
Call $\varepsilon=0$ the \emph{case choice} and $\varepsilon=1$
the \emph{progress choice}. Both retain exactly one ordinary
formula assumption. Compound matrices have only the case choice.

In DI the eigenparameter $z$ is fresh for $\Gamma$, the witness
$t$, and every fixed parameter of $B$. In DE the eigenparameter
$a$ is fresh for $\Gamma$, the conclusion $F$, and every fixed
parameter of $B$. In transport, $F(v)$ is any fixed native
formula template. Displayed local assumptions and clauses are
discharged. There are no primitive reflexivity, coverage,
separation, ordinary logical, or additional induction rules.

We call this calculus $\mathcal C$. The original comparison
calculus $\mathcal N$ has identical syntax, computation,
Use, DE, and transport, but its introduction is
\[
\frac{\Gamma,B(t,0),[B(t,z)]_{\mathrm{at}}\vdash B(t,Sz)}
     {\Gamma\vdash\mathsf D_{x,y}B(x,y)}.
\tag{Old DI}\label{eq:choice-old-di}
\]
The bracket supplies $B(t,z)$ exactly at a literal equation.

\subsection{Exact equivalence and the unused-witness construction}

\begin{lemma}[Structural admissibility]
\label{lem:choice-structural}
Weakening, exchange, contraction, renaming, ordinary-assumption cut,
and capture-free numerical substitution are admissible in $\mathcal C$.
\end{lemma}
\begin{proof}
Induct on derivations. Rename local eigenparameters before inserting
proofs or terms. Numerical substitution preserves a formula's root
constructor, so an allowed progress choice remains allowed.
The selected choice does not change under substitution or cut.
For ordinary-assumption cut, replace each use of its constant clause
by the supplied proof. Computation schemas are closed under
capture-free numerical substitution. These facts verify every case.
\end{proof}

\begin{theorem}[Identity-language equivalence]
\label{thm:choice-equivalence}
For every retained clause context $\Gamma$ and native formula $F$,
\[
 \Gamma\vdash_{\mathcal C}F
 \quad\Longleftrightarrow\quad
 \Gamma\vdash_{\mathcal N}F.
\]
The transformations are effective on finite derivations and leave
all numerical and formula syntax unchanged.
\end{theorem}
\begin{proof}
Each new introduction is an instance of the old one after weakening
by the omitted local assumption. For the case choice at an equation,
add the induction instance. For the progress choice, add the zero
instance. Compound case introduction is already the old rule.
Thus every new derivation translates to an old one.

For the reverse direction, transform an old derivation recursively.
Only its atomic DI case needs a construction. Suppose its matrix is
the equation $E(x,y)$ and its transformed premise is
\[
 \pi:\quad\Gamma,E(t,0),E(t,z)\vdash E(t,Sz).
\]
The outer new inference will be case DI with witness $t$. It is
therefore enough to prove $E(t,Sz)$ under
$\Gamma,E(t,0)$.

Rename the eigenvariable in $\pi$ to a fresh $v$. Choose a new
variable $u$ absent from $E(t,y)$ and all retained data. Under
$\Gamma,E(t,0)$, progress DI with witness zero introduces
\[
             \mathsf D_{u,y}E(t,y)
\]
from exactly the renamed premise $\pi[v/z]$.
The bound coordinate $u$ is deliberately unused.

Apply DE to this paired formula. Because $u$ is unused, the opened
clause is exactly
\[
               \langle E(t,0);y.E(t,y)\rangle;
\]
it contains no new existential witness. The retained ordinary
assumption $E(t,0)$ supplies its antecedent. Use at $Sz$ gives
$E(t,Sz)$, and DE closes. Finally the outer case DI introduces
the original $\mathsf D_{x,y}E(x,y)$ with witness $t$.
The original premise proof $\pi$ has been inserted once.

An old compound introduction already is the new case choice.
All other rule cases are unchanged. This completes the induction
on the original derivation.
\end{proof}

The unused inner coordinate is essential. Introducing
$\mathsf D_{x,y}E(x,y)$ internally and then eliminating it would
give an arbitrary fresh witness, not necessarily the original
$t$. The construction instead keeps $t$ as a fixed parameter
of $E(t,y)$ and binds only the unused coordinate $u$.

\begin{proposition}[Linear bound on inference nodes]
\label{prop:choice-node-bound}
Count each computation axiom, Use, DI, DE, and transport as one
inference node, expanding ordinary-assumption use as a computation
of $H$ followed by Use. If an old derivation has $N$ nodes and
$r$ atomic DI nodes, the displayed transformation has exactly
\[
                   N+5r\leq 6N
\]
nodes.
\end{proposition}
\begin{proof}
Each atomic old DI node is replaced by six nodes: outer case DI,
inner progress DI, DE, Use of the opened clause, Use of the
retained base assumption, and the computation of $H$ needed for
that last ordinary-assumption use. Its old premise subtree
occurs once. All other inference nodes are copied once.
\end{proof}

This counts inference-tree nodes. It is not a bound on serialized
JSON bytes, formula sizes, checker runtime, or the number of
consumed packet cells when hypothesis closure is implicit.
In the formula-context Q/D presentation, the corresponding
replacement has outer case DI, inner progress DI, DE, QE,
and two Hyp nodes, giving the same six-node count.

\subsection{Induction and numerical cases separate cleanly}

\begin{lemma}[Contextual equation induction]
\label{lem:choice-ei}
The progress choice, DE, and Use derive
\[
\frac{\Gamma\vdash E(0)\qquad
      \Gamma,E(z)\vdash E(Sz)}
     {\Gamma\vdash E(t)}
\tag{EI}\label{eq:choice-ei}
\]
for every equation $E$, with the usual freshness condition.
\end{lemma}
\begin{proof}
Choose an unused witness coordinate $u$. Progress DI introduces
$\mathsf D_{u,y}E(y)$ from the given step. DE opens the clause
$\langle E(0);y.E(y)\rangle$. Supply its antecedent using the given
base proof and use the clause at $t$. Close DE.
\end{proof}

\begin{lemma}[Cases, equality, and binder congruence]
\label{lem:choice-preinduction}
The case choice derives, for arbitrary native formulas,
\[
\frac{\Gamma\vdash F(0)\qquad\Gamma\vdash F(Sz)}
     {\Gamma\vdash F(t)}
\tag{NC}\label{eq:choice-nc}
\]
with the usual freshness condition. Equality reflexivity,
symmetry, transitivity, numerical congruence, and uniform
matrix interderivability under $\mathsf D$ are derivable
without using the progress choice.
\end{lemma}
\begin{proof}
For cases, use case DI with an unused witness coordinate and
matrix $F(y)$, ignoring its retained base in the successor
premise. DE opens $\langle F(0);y.F(y)\rangle$; the given zero
proof and Use at $t$ finish.

For reflexivity put $r=I(0,t;u.0)$. Use the computation equation
$r=t$ as both transport premises with motive $v=t$ to obtain
$t=t$. The usual explicit equality motives give symmetry,
transitivity, and numerical congruence.

For binder congruence, open the old paired formula at a fresh
witness $a$ and introduce the new one by case DI at $a$.
Convert its retained base to the old base, use the opened
old clause at the successor, and convert back. Ordinary cut
closes the construction, and the reverse conversion gives
the converse. No induction is used.
\end{proof}

The primitive progress choice is still restricted to equations.
Allowing it at arbitrary compound matrices would, by the
argument for \eqref{eq:choice-ei}, supply induction for
arbitrary formulas and change the intended theory.

\subsection{The entire ordinary logical interface}
\label{subsec:choice-logic}

Here is a predecessor-free reconstruction of the ordinary logical
operations. It makes explicit that the case choice supplies the
logical interface independently of progress induction.

Let $C(v,a,b):=I(v,a;u.b)$ with $u$ unused. Computation gives
$C(0,a,b)=a$ and $C(Sv,a,b)=b$. Define the gate on fixed finite
formula syntax by
\[
 G_v(s=t):=(C(v,0,s)=C(v,0,t)),\qquad
 G_v(\mathsf D_{x,y}B):=\mathsf D_{x,y}G_vB,
\]
renaming bound variables away from the control. Equality,
binder congruence, and external structural induction give
\[
 \vdash G_0A,\qquad G_{Sv}A\dashv\vdash A,\qquad
 A\dashv\vdash\mathsf D_{u,y}G_yA
 \quad(u,y\notin\operatorname{FV}(A)).
\tag{Gate and padding}\label{eq:choice-gate}
\]
For padding, case DI proves the forward direction. DE, the
zero-gate theorem, and Use at $1$ prove the reverse direction.

Define a selector between two arbitrary formulas by
\[
 M_v(s=t,u=w):=(C(v,s,u)=C(v,t,w))
\]
and, with aligned fresh paired variables, by
\[
\begin{aligned}
 M_v(\mathsf D_{x,y}A,\mathsf D_{x,y}B)
   &:=\mathsf D_{x,y}M_v(A,B),\\
 M_v(E,\mathsf D_{x,y}B)
   &:=\mathsf D_{x,y}M_v(G_yE,B),\\
 M_v(\mathsf D_{x,y}A,E)
   &:=\mathsf D_{x,y}M_v(A,G_yE),
\end{aligned}
\]
where $E$ is an equation. Every call decreases maximum input
binder depth. Padding and congruence prove
$M_0(A,B)\dashv\vdash A$ and
$M_{Sv}(A,B)\dashv\vdash B$.

Put $\delta(v)=C(v,1,0)$ and $p(z)=I(z,0;u.\delta(u))$.
Only the fixed computations $p(1)=1$ and $p(2)=0$ are needed.
The full logical macro inventory is
\[
\begin{aligned}
 A\Rightarrow_*B&:=\mathsf D_{u,z}M_z(A,B),\\
 A\wedge_*B&:=\mathsf D_{u,z}G_zM_{p(z)}(A,B),\\
 \exists_*x\,A(x)&:=\mathsf D_{x,y}G_yA(x),\\
 A\vee_*B&:=\exists_*x\,M_x(A,B),\\
 \forall_*x\,A(x)&:=
        A(0)\wedge_*\mathsf D_{u,z}G_zA(z),\\
 \bot_*&:=(0=1).
\end{aligned}
\tag{Logical macros}\label{eq:choice-macros}
\]
Newly bound variables are fresh; $u$ is unused.

Implication introduction obtains its assumption from the zero
selector and supplies the consequent at successors, then uses
case DI. Elimination opens its clause and uses it at $1$.
Existential introduction uses its witness and the successor
gate; elimination opens at a fresh witness, supplies the
zero-gate theorem, and uses the clause at $1$.

Disjunction introduction uses witnesses $0$ and $1$.
For elimination, discharge the two branch proofs as
$A\Rightarrow_*F$ and $B\Rightarrow_*F$. NC applied to
$M_z(A,B)\Rightarrow_*F$ gives this implication at every
control: the two endpoint cases follow from the branch
implications. Existential elimination provides $M_a(A,B)$
and the implication gives $F$.

Given $A$ and $B$, selector endpoints and NC give $M_v(A,B)$
at an arbitrary control. Substitute $v=p(Sz)$ and use the
successor gate for conjunction introduction. For projections,
open its zero-gated clause and specialize at $1$ and $2$;
the two fixed computations give $B$ and $A$, respectively.

These rules prove existential numerical coverage:
\[
     (x=0)\vee_*\exists_*v\,(x=Sv).
\tag{Coverage}\label{eq:choice-coverage}
\]
Use NC. Zero takes the left injection of reflexivity, and a
successor takes the right injection with the evident witness.
Coverage gives equality cases even when the retained context
contains $x$; its elimination adds $x=0$ or $x=Sv$ rather
than silently changing that retained context.

Universal introduction follows from a fresh-parameter proof
by substituting zero and successors and forming the two
conjuncts. For elimination, take its $A(0)$ conjunct and
open the paired conjunct. Its base is a zero-gate theorem.
Coverage of the target chooses the zero instance from
$A(0)$ or a successor instance from the opened clause
and gate removal.

Finally, under $0=1$, computation and transport give
$s=C(0,s,t)=C(1,s,t)=t$ for arbitrary numerical terms.
Every equation follows, and external induction on formula
syntax with case DI gives every native formula. This
supplies falsity elimination. No excluded middle for
quantified formulas has entered the construction.

\subsection{Whole-domain realizations and both formula composites}

For a standard first-order reading, let $T_{\rm eq}$ have
intuitionistic equality logic, the same numerical signature
and computation equations, numerical coverage, separation
$(0=1)\to\bot$, and contextual equation induction.
Its interpretation of the native syntax is
\[
 (s=t)^\circ=(s=t),\qquad
 (\mathsf D_{x,y}B)^\circ
   =\exists x\bigl(B^\circ(x,0)\to\forall y B^\circ(x,y)\bigr),
\]
\[
 \langle A;y.B\rangle^\circ
       =A^\circ\to\forall y B^\circ.
\]
The domain is the entire numerical domain and interpreted
equality is actual equality.

Case DI is sound in this realization by numerical coverage.
Progress DI is sound by contextual equation induction under
the conditional base. Use, DE, transport, and computation
have their displayed ordinary meanings. Conversely, the
native macros derive every ordinary logical rule,
\eqref{eq:choice-ei} derives equation induction, and
\eqref{eq:choice-coverage} derives coverage. Separation
translates to the identity implication
$(0=1)\Rightarrow_*(0=1)$.

For the essential paired formula composite,
\[
 \mathsf D_{x,y}B
 \dashv\vdash
 \exists_*x\bigl(B(x,0)\Rightarrow_*\forall_*y\,B(x,y)\bigr).
\tag{Paired composite}\label{eq:choice-paired-composite}
\]
Forwards, DE opens its clause at a fresh $a$; under $B(a,0)$
that clause supplies $B(a,y)$ at a fresh arbitrary $y$.
Introduce the universal, implication, and existential
macros and close DE. Backwards, eliminate the existential
at a fresh $a$. Introduce the paired formula with witness
$a$ by case DI: the retained base, supplied implication,
and universal elimination at $Sz$ give its successor premise.

The ordinary interpretations of the six logical macros
have their ordinary round trips. For conjunction, its true
base and the instances $1,2$ recover both operands;
the converse uses selector cases. The universal has all
successor instances in its paired conjunct and zero in
its first conjunct; coverage gives exactly the ordinary
universal. Implication, disjunction, and existential
quantification use their selector or gate endpoints.
Falsity uses the explicit ordinary separation bridge.
External induction on finite formulas, together with
\eqref{eq:choice-paired-composite}, gives both formula
composites. Clauses correspond in both directions to
$A\Rightarrow_*\forall_*yB$ by the derived logical rules.

The numerical composites are literal identities with
internally provable graph $x=y$. No standard-numeral
restriction, retract, quotient, or equality reflection
assumption is being inserted at this stage. The separately
supplied arithmetic bridge identifies $T_{\rm eq}$ with
full iPRA; the present change preserves that bridge
unchanged because Theorem~\ref{thm:choice-equivalence}
preserves every native sequent.

\subsection{What the theorem does and does not claim}

The new inference interface retains one ordinary assumption
at every paired introduction. It preserves the entire
consequence relation, not just equations or closed
theorems. Consequently the established full-iPRA
interpretations and equiconsistency transfer unchanged.

This result is different from retaining only progress
introduction and strengthening closure elimination.
Here the original elimination rule is retained, and
the case and progress introductions coexist. It is also
different from allowing arbitrary compound progress,
which would supply stronger induction.

A complete-checker comparison must still count the
certificate choice, its equation-root restriction,
substitution, context handling, computation, and the
public input/output interface. The mathematical theorem
alone is not a byte-count or minimality theorem.
The inference-node bound is distinct from source-size
measurements and from serialized certificate size.
The proof is a written universal proof; finite regression
tests do not constitute a Lean/Rocq formalization of it.

% End included source


% Begin included source: reference/inherited_ordinary_bridge.tex
\section{Arithmetic strength and both ordinary interpretation composites}
\label{sec:ordinary-bridge}

We spell out the arithmetic bridge to avoid identifying standard numerical
agreement with bi-interpretability. First elaborate every fixed iterator
into an ordinary function symbol. At stage zero the symbols are $0,S$.
Given a fixed earlier-stage term $h(u,\vec p)$, introduce
$i_h(n,a,\vec p)$ with the two equations
\[
i_h(0,a,\vec p)=a,\qquad
i_h(Sn,a,\vec p)=h(i_h(n,a,\vec p),\vec p).
\]
Only finite descriptions occur; the domain contains no code for a
universal evaluator. Formal parameters are distinct when a description
is declared, while actual arguments can coincide.

Let $T_I$ be ordinary intuitionistic first-order logic with actual
equality and genuine falsity, these staged equations, contextual
induction for literal equations, and
\[
\forall x\,(x=0\vee\exists y\,x=Sy),
\qquad (0=S0)\to\bot.
\]
The second sentence is the falsity bridge. It must be stated: positive
numerical equations by themselves admit a singleton structure. Native
explosion from $0=S0$ and a translation of ordinary falsity to that
equation do not authorize silently dropping the bridge from $T_I$.

The native-to-ordinary translation is
\[
(s=t)^\circ=(s=t),\qquad
(\mathsf D_{x,y}B)^\circ=
\exists x\bigl(B^\circ(x,0)\to\forall y B^\circ(x,y)\bigr).
\]
A clause $\langle A;y.B\rangle$ becomes $A^\circ\to\forall yB^\circ$.
For the auxiliary closure used by the compact checkers,
$(\mathsf Q_yB)^\circ=B^\circ(0)\to\forall yB^\circ$.
The interpretation domain is all elements, defined by $x=x$;
equality is actual equality.

Case DI is sound because numerical coverage combines its retained base
with its successor proof. Progress DI is sound by literal-equation
induction under a temporarily assumed base. Use, DE and transport are
ordinary logical inferences under the translations. Conversely the
derived logical operations, equation induction, cases and explosion in
Section~\ref{sec:choice-introduction} recover the rules and axioms of
$T_I$. An auxiliary closure is interderivable with
$\mathsf D_{v,y}B(y)$ for unused $v$: in one direction use case DI with
the closure and its base; in the other eliminate the paired formula
and return its supplied closure. Thus formulas and context clauses have
the stated meanings in both presentations.

\subsection{The arithmetic dependency order}

The scalar update has no recursion-counter argument. We do not conceal
one in an encoding. The inherited internal proofs, reproduced in the
appendices, establish the required arithmetic in this order:
\begin{enumerate}
\item Using equation induction and numerical cases, construct a fixed
scalar predecessor term and prove $P0=0$ and $P(Sn)=n$.
Existential coverage then gives $x=0\vee x=S(Px)$: in the successor
case $x=Sy$, apply $P$ to obtain $Px=y$ and substitute.
\item Construct truncated subtraction and a numerical characteristic
of actual equality. Prove its reflection on the entire original domain.
\item Compile each fixed quantifier-free formula into its characteristic
equation and recover every contextual quantifier-free induction instance.
\item With that induction available, construct pairing and projections
by a bounded square-root search and compile every fixed general
primitive-recursive description into a scalar iterator term.
\end{enumerate}
No pairing or faithful equality detector is assumed to obtain
predecessor. No appeal to standard-model correctness replaces an
internal equality. Every induction fixes its passive parameters;
diagonal substitution occurs only into a previously proved identity.

For a target recursion
\[
f(0,\vec p)=g(\vec p),\qquad
f(Sn,\vec p)=h(n,f(n,\vec p),\vec p),
\]
let $\pi,L,H$ be the proved pairing and projections. The compiler uses
\[
V(n,\vec p)=I\bigl(n,\pi(0,g^*(\vec p));
u.\pi(SL(u),h^*(L(u),H(u),\vec p))\bigr),
\qquad f^*(n,\vec p)=H(V(n,\vec p)).
\]
Equation induction proves $L(V(n,\vec p))=n$; the two projection
identities then prove the target defining equations. The construction
is external recursion on each finite PR description, not an internal
quantification over arbitrary functions.

\begin{theorem}[Whole-domain bi-interpretability and equiconsistency]
The ordinary realization $T_I$ and full first-order iPRA are
bi-interpretable. In both theories $x=y$ defines the isomorphism from
the original domain to its double interpretation. The native Choice-DI
calculus is therefore a faithful proof presentation of a theory
bi-interpretable and equiconsistent with iPRA.
\end{theorem}
\begin{proof}
Interpret an iterator in iPRA by primitive recursion that ignores its
counter argument in the update. iPRA proves its equations, literal
equation induction, coverage and separation. In the reverse direction
use the scalar compiler. The predecessor and equality-reflection proofs
supply the successor axioms and every quantifier-free induction instance.
All ordinary logical rules are preserved.

For two iterators with the same seed and provably equal updates,
equation induction on the equality of their values proves uniqueness at
every counter, with all parameters fixed. The corresponding argument
applies to ordinary primitive recursion. Induct externally on a finite
description: its double translation satisfies its original defining
equations, so internal uniqueness identifies the two functions. This
also proves coherence when substitution changes the staged syntactic
description of an update.

The domain and equality are unchanged; hence the identity graph is
internally total, bijective and preserves every interpreted operation.
The displayed ordinary meanings of the logical macros and the paired
binder give both formula composites by external induction on finite
formula syntax. Clauses are similarly equivalent to their conditional
universal meanings. The explicit falsity bridge preserves contradiction
in both directions. Effective translations of finite proofs give
equiconsistency, without a claim that either theory proves its own
consistency.
\end{proof}

% End included source


% Begin included source: reference/scalar_predecessor.tex
% Adapted with attribution from the pinned unparenthesized one-hole archive.
\section{Internal scalar predecessor bootstrap}
\label{sec:hybrid-bootstrap}

This scalar bootstrap is inherited from the supplied scalar development~\cite{erasure} and is reproduced to keep the scalar checker self-contained. It is not a new contribution of this paper. The proof uses a factor map fixing odd states and a saturated-iteration construction of predecessor.

The available arithmetic terms are $0$, $S$, and
\[
 I(n,a;u.h),\qquad
 I(0,a;u.h)=a,\qquad
 I(Sn,a;u.h)=h[I(n,a;u.h)/u].
\]
The update binds only $u$.  All other variables in it are fixed numerical
parameters.  We use equality reasoning, the already reconstructed
zero/successor cases, and induction on a single literal equation.
Every induction below keeps its other parameters fixed.  Predecessor,
successor injectivity, pairing, and induction on arbitrary
quantifier-free formulas are not assumptions of the argument.

\subsection{Elementary terms and two finite cycles}

Write $1=S0$ and use larger numerals as successor abbreviations.  Define
\[
 C(n,a,b)=I(n,a;u.b),\qquad
 Z(n)=C(n,1,0),\qquad
 a+b=I(b,a;u.Su).
\]
The update variable of $C$ is absent from its body.  Consequently
\[
 C(0,a,b)=a,\qquad C(Sn,a,b)=b.
\]
Numerical cases show that $Z(n)$ is $0$ or $1$.
The usual commutative-monoid equations for addition follow by literal
equation induction.  For example, $0+b=b$ and $Sa+b=S(a+b)$ are
proved by induction on $b$; these yield commutativity by a further
equation induction, and associativity follows by induction on the last
summand.  We use these equations to reorder finite sums.  The notation
$2a$ means $a+a$ and $3a$ means $a+a+a$; no multiplication operation
is needed for this bootstrap.

Put
\[
 N(n)=I(n,0;u.Z(u)),\qquad
 j(u)=Z(u)+2N(u),\qquad
 M(n)=I(n,0;u.j(u)).
\]
The range of $N$ is contained in $\{0,1\}$: its zero value is $0$,
and its successor value is a value of $Z$.  This uses numerical cases
on its counter, not induction on a disjunction.  Hence
\begin{equation}
 N(n)+Z(N(n))=1,\qquad
 Z(Z(N(n)))=N(n),\qquad N(n+2)=N(n).
 \label{eq:hybrid-parity}
\end{equation}
Here the last equation follows from two successor computations and
the preceding equation.  Literal-equation inductions give
\begin{equation}
 N(2n)=0,\qquad N(2n+1)=1,\qquad N(3n)=N(n).
 \label{eq:hybrid-parity-scaling}
\end{equation}
For the first equation the successor step adds two to the argument of
$N$; the second follows by one computation.  In the third, increasing
$n$ adds three to the left argument and therefore applies $Z$ once to
both sides of the induction equation.

Finite computations now give
\[
 j(0)=1,\qquad j(1)=2,\qquad j(2)=0.
\]
The proof below needs just these three values and
$M(Sn)=j(M(n))$; it does not assume a division algorithm or a theorem
characterizing $M$ as remainder modulo three.

\subsection{A ramp with a global scaling identity}

Define
\[
 G(n)=I(n,0;u.u+M(u)+1).
\]
The central invariant is
\begin{equation}
 M(G(n))=N(n).
 \label{eq:hybrid-ramp-invariant}
\end{equation}
Its zero case is immediate.  For its successor step, retain the
single induction equation and split the Boolean value $N(n)$.
If it is $0$, the next $G$ state is $G(n)+1$; its $M$ value is
$j(0)=1$, equal to $N(Sn)$.  If it is $1$, the next $G$ state
is $G(n)+2$; two computations of $M$ change its value from $1$
to $2$ and then to $0$, again equal to $N(Sn)$.  This proves
the invariant by literal-equation induction.

It follows that
\[
 G(Sn)=G(n)+N(n)+1,\qquad G(n+2)=G(n)+3.
\]
Equation induction therefore yields
\begin{equation}
 G(2n)=3n.
 \label{eq:hybrid-ramp-double}
\end{equation}
A separate literal-equation induction proves
\begin{equation}
 G(3n)=3G(n)+N(n).
 \label{eq:hybrid-ramp-triple}
\end{equation}
For completeness, its successor calculation is
\begin{align*}
 G(3Sn)
 &=G(3n)+N(3n)+4
   =3G(n)+2N(n)+4,\\
 3G(Sn)+N(Sn)
 &=3G(n)+3N(n)+3+Z(N(n)).
\end{align*}
The expressions agree in both Boolean cases for $N(n)$.
The base is $0=0$.  In particular,
\eqref{eq:hybrid-ramp-triple} was proved internally; it was not
inferred from an intended floor-function interpretation of $G$.

Set
\[
 H(u)=C(N(u),G(u),u).
\]
The crucial choice is to retain $u$ when its parity is $1$.
The preceding equations give
\begin{equation}
 H(2u)=3u,\qquad H(1)=1,\qquad H(3u)=3H(u).
 \label{eq:hybrid-factor-map}
\end{equation}
To check the last identity, use $N(3u)=N(u)$ and split the Boolean
$N(u)$.  At $0$, both sides are $3G(u)$ by
\eqref{eq:hybrid-ramp-triple}; at $1$, both are $3u$.
This argument does not require that an arbitrary number has already
been represented as $2v$ or $2v+1$.

\subsection{A comparator from scalar iteration}

For every fixed update $h$, equation induction on $n$, keeping $m,a$
and all parameters fixed, proves the semigroup identity
\[
 I(m+n,a;h)=I(n,I(m,a;h);h).
\]
The base is the zero computation and the step applies $h$ to the
induction equation. This identity will also be used in the later
pairing construction.


For a fixed update $h$, equation induction on $y$, with its initializer
$a$ and all parameters of $h$ fixed, proves the rotation identity
\begin{equation}
 I(Sy,a;h)=I(y,h(a);h).
 \label{eq:hybrid-rotation}
\end{equation}
Both sides have base $h(a)$, and their successor equations apply the
same $h$ to the previous equality.  Here and below $I(y,a;h)$ is
an abbreviation for one fixed finite scalar term, not quantification
over functions.

By \eqref{eq:hybrid-factor-map}, another equation induction gives
\begin{equation}
 I(y,3a;H)=3I(y,a;H).
 \label{eq:hybrid-iterator-scaling}
\end{equation}
Its base is $3a=3a$.  In the step the fixed-parameter induction
equation rewrites the new left side as
\[
 H(3I(y,a;H))=3H(I(y,a;H)).
\]
The initializer parameter stays $a$ throughout this induction.
Equation induction also gives $I(y,1;H)=1$, since $H(1)=1$.

Define
\[
 E(x)=I(x,1;u.2u),\qquad
 L(x,y)=N\bigl(I(y,E(x);H)\bigr).
\]
Then the following entire recursion specification is provable:
\begin{equation}
 L(0,y)=1,\qquad L(Sx,0)=0,\qquad L(Sx,Sy)=L(x,y).
 \label{eq:hybrid-comparator}
\end{equation}
The first equation uses $E(0)=1$ and the fixed-point identity for
$H$.  The second uses $E(Sx)=2E(x)$ and $N(2a)=0$.
For the third,
\begin{align*}
 L(Sx,Sy)
 &=N\bigl(I(Sy,2E(x);H)\bigr)\\
 &=N\bigl(I(y,H(2E(x));H)\bigr)\\
 &=N\bigl(I(y,3E(x);H)\bigr)\\
 &=N\bigl(3I(y,E(x);H)\bigr)\\
 &=N\bigl(I(y,E(x);H)\bigr)=L(x,y).
\end{align*}
The successive steps use rotation, $H(2a)=3a$, iterator scaling,
and $N(3a)=N(a)$.  Thus $L$ has the comparison equations at all
arguments internally.  Its standard-model meaning as the indicator
of $x\le y$ is motivation, not a premise of the following argument.

\subsection{Saturated iteration gives predecessor}

We now need only the three equations
\eqref{eq:hybrid-comparator}.  Define
\[
 B_m(u)=C(L(Su,m),u,Su),\qquad
 W(t,m)=I(t,0;u.B_m(u)).
\]
The subscript in $B_m$ records a fixed parameter; $B_m$ is an open
term abbreviation, not a function variable.

Numerical cases on $e$ prove the selector equation
\[
 C(e,Sa,Sb)=S C(e,a,b).
\]
Equivalently one may use equation induction on $e$, whose induction
hypothesis is unused.  Therefore
\begin{align}
 B_{Sm}(Su)
 &=C(L(SSu,Sm),Su,SSu)\notag\\
 &=C(L(Su,m),Su,SSu)=S B_m(u),
 \label{eq:hybrid-conjugacy}\\
 B_{Sm}(0)&=C(L(S0,Sm),0,S0)=C(L(0,m),0,S0)=S0.
 \label{eq:hybrid-conjugacy-base}
\end{align}

Literal-equation induction on $t$, with $m$ fixed, proves
\begin{equation}
 W(St,Sm)=S W(t,m).
 \label{eq:hybrid-saturated-shift}
\end{equation}
The base is \eqref{eq:hybrid-conjugacy-base}.  The successor step is
\begin{align*}
 W(SSt,Sm)
 &=B_{Sm}(W(St,Sm))\\
 &=B_{Sm}(S W(t,m))\\
 &=S B_m(W(t,m))=S W(St,m).
\end{align*}
Only the second line uses the current induction hypothesis, at the
same fixed $m$.

Two subsequent equation inductions give
\begin{equation}
 L(Sn,n)=0,\qquad W(n,n)=n.
 \label{eq:hybrid-diagonals}
\end{equation}
The first uses the comparator recurrence.  The second uses
\eqref{eq:hybrid-saturated-shift} in its successor step.  The
substitution $t=m=n$ is made into that already proved universally
parameterized equation; no parameter of an active induction hypothesis
is changed.  Consequently
\begin{equation}
 W(Sn,n)=B_n(n)=C(L(Sn,n),n,Sn)=n.
 \label{eq:hybrid-extra-step}
\end{equation}

Finally define the finite scalar term
\begin{equation}
 P(m)=I\bigl(m,0;u.C(L(SSu,m),u,Su)\bigr).
 \label{eq:hybrid-predecessor-term}
\end{equation}
Its zero equation is immediate.  To justify its successor equation
without assuming a separate extensionality rule for iterator bodies,
put
\[
 V(t,m)=I\bigl(t,0;u.C(L(SSu,Sm),u,Su)\bigr).
\]
The comparator recurrence identifies this update pointwise with
$B_m(u)$.  Equation induction on $t$, with $m$ fixed, gives
$V(t,m)=W(t,m)$: their base values are $0$, and the successor step
uses the pointwise update equation and then the induction equation.
Now \eqref{eq:hybrid-extra-step} yields
\[
 \boxed{\qquad P(0)=0,\qquad P(Sm)=V(Sm,m)=W(Sm,m)=m.\qquad}
\]

\subsection{Dependency boundary}

All terms in this section are built from $0,S,I$ by finite substitution.
All inductions have a single equation as their invariant.  Numerical
cases and equality transport are the already reconstructed logical
resources.  Successor injectivity may now be obtained by applying $P$
to an assumed equation $Sx=Sy$.  Equality tests, reflection of
quantifier-free formulas, pairing, and the recovery of general
primitive recursion can then be developed in the established order.
They were not used in order to manufacture predecessor.

This argument is a written internal derivation in the stated theory.
The finite implementation checks do not constitute a proof-assistant
formalization of this entire metatheorem.

% End included source


% Begin included source: reference/equality_reflection.tex
% Standalone section.  Requires only amsmath and amssymb in the parent document.
% No theorem environments, macros, or bibliography keys are assumed.
% All arithmetic claims in this section are proved here.

\section{From literal-equation induction to full intuitionistic PRA}
\label{sec:eic-completion}\label{app:arithmetic}

This section supplies an unconditional completion theorem. It is adapted
with attribution from the current forward-left draft~\cite{forward}. In the
scalar application used here, every function occurring in the proof is
already a scalar term: predecessor is supplied by
Appendix~\ref{sec:hybrid-bootstrap}, the selector is
$C(n,a,b)=I(n,a;u.b)$, addition is $x+y=I(y,x;u.Su)$, and
truncated subtraction is $d(x,y)=I(y,x;u.Pu)$. All other functions below
are finite compositions of these. Thus this proof does not presuppose
that arbitrary primitive recursion has already been compiled. The QFI
conclusion first applies to the scalar language; the general PR compiler
comes afterward.  Literal-equation
induction, together with the primitive-recursive defining equations,
predecessor coverage, and separation of $0$ and $1$, proves induction for every
quantifier-free formula.  In particular, the proof does not assume a faithful
numerical equality test, stability or decidability of equality, successor
uniqueness, or any commutation law for an iterator.  We construct the equality
test and prove its reflection property for arbitrary elements of the
numerical domain.

\subsection{The exact hypotheses}

Work in ordinary intuitionistic first-order logic with equality and genuine
falsity $\bot$.  The language has one numerical sort and a symbol for each
fixed, finite primitive-recursive description.  Distinct descriptions are
not identified merely because they compute the same standard function.
Let $\mathsf{PRE}$ contain the defining equations for zero, successor,
projections, composition, and primitive recursion.  In particular, choose
descriptions $P$ and $C$ with equations
\begin{align}
 P(0)&=0, & P(Sx)&=x, \label{eq:eic-predecessor}\\
 C(0,a,b)&=a, & C(Sx,a,b)&=b. \label{eq:eic-selector}
\end{align}
Write $1=S0$.  In addition to $\mathsf{PRE}$, assume
\begin{align}
 \mathsf{Cov}:&\qquad
       \forall x\,\bigl(x=0\lor x=S(Px)\bigr),
       \label{eq:eic-coverage}\\
 \mathsf{Sep}:&\qquad 0=1\longrightarrow\bot,
       \label{eq:eic-separation}
\end{align}
and induction for each \emph{literal equation} $E(n,\vec p)$:
\begin{equation}
 \forall\vec p\,
 \left[
   \left(E(0,\vec p)\land
       \forall n\bigl(E(n,\vec p)\longrightarrow E(Sn,\vec p)\bigr)
   \right)
   \longrightarrow\forall n\,E(n,\vec p)
 \right].
 \label{eq:eic-induction}
\end{equation}
Here the schematic expression $E$ is a single equality of numerical
terms; the displayed universal quantifiers range only over numbers.
Every description and every instance of the schema is finite external
syntax.

The axiom formulation \eqref{eq:eic-induction} and the usual contextual
induction rule are interchangeable.  For example, if
\[
 \Gamma\vdash E(0,\vec p),
 \qquad
 \Gamma,E(z,\vec p)\vdash E(Sz,\vec p),
\]
where $z$ is absent from $\Gamma$ and the passive parameters, discharge the
second premise's equation and universally introduce $z$.  The axiom then
gives $\Gamma\vdash\forall n E(n,\vec p)$, which can be instantiated at any
target term.  Conversely, assume the antecedent of the axiom and apply the
contextual rule with fresh variables.  This uses the ordinary intuitionistic
logical rules and permits arbitrary, including quantified, formulas in
$\Gamma$.

\subsection{Constructor separation and two elementary tests}

Define the sign description and the countdown description by
\begin{equation}
 \operatorname{sg}(u)=C(u,0,1),
 \qquad
 K(u)=C(u,1,Pu).
 \label{eq:eic-sign-countdown}
\end{equation}
Their defining equations imply
\begin{align}
 \operatorname{sg}(0)&=0,
 &\operatorname{sg}(Sx)&=1,\label{eq:eic-sign-computation}\\
 K(0)&=1,
 &K(Sx)&=x.\label{eq:eic-countdown-computation}
\end{align}

\paragraph{Constructor laws.}
Successor is injective: applying $P$ to $Sx=Sy$ gives $x=y$ by
\eqref{eq:eic-predecessor}.  Successor never equals zero.  Indeed,
$Sx=0$ gives $1=0$ by applying $\operatorname{sg}$ and using
\eqref{eq:eic-sign-computation}; symmetry and
\eqref{eq:eic-separation} give $\bot$.  Consequently $1\ne0$, where
$A\ne B$ abbreviates $A=B\to\bot$.  None of these arguments uses
induction.

\paragraph{The sign test.}
For every $u$,
\begin{align}
 &\operatorname{sg}(u)=0\lor\operatorname{sg}(u)=1,
       \label{eq:eic-sign-range}\\
 &\operatorname{sg}(u)=0\quad\longleftrightarrow\quad u=0,
       \label{eq:eic-sign-zero}\\
 &\operatorname{sg}(u)=1\quad\longleftrightarrow\quad u\ne0.
       \label{eq:eic-sign-one}
\end{align}
To prove the range assertion, split \eqref{eq:eic-coverage}.  The branch
$u=0$ gives sign $0$, and the branch $u=S(Pu)$ gives sign $1$.
The implication $u=0\to\operatorname{sg}(u)=0$ is computation.
For its converse, the successor branch of coverage would make
$\operatorname{sg}(u)=0$ imply $1=0$, a contradiction; in that branch
explosion gives $u=0$ as well.  Similarly, if
$\operatorname{sg}(u)=1$ and $u=0$, computation gives $0=1$.
Conversely, $u\ne0$ excludes the zero branch of coverage, and the
successor branch computes the sign as $1$.

Coverage and \eqref{eq:eic-selector} also give, for all $u,v$,
\begin{equation}
 C(u,v,v)=v.
 \label{eq:eic-selector-idempotence}
\end{equation}
Both coverage branches compute to $v$, so this identity again needs no
induction.

\paragraph{The countdown test and the notation $W$.}
For every $u$,
\begin{equation}
 K(u)=0\quad\longleftrightarrow\quad u=1.
 \label{eq:eic-countdown-zero}
\end{equation}
If $u=1=S0$, \eqref{eq:eic-countdown-computation} gives $K(u)=0$.
Conversely, assume $K(u)=0$ and split coverage.  If $u=0$, then
$K(u)=1$, contradicting separation, and explosion gives $u=1$.
If $u=S(Pu)$, computation gives $K(u)=Pu$; hence $Pu=0$ and
$u=S0=1$.

If a presentation uses a description named $W$ with defining equations
\[
 W(0)=1,\qquad W(Sx)=x,
\]
then it proves $K(u)=W(u)$ for every $u$ already from coverage.
In the zero branch both sides are $1$; in the branch $u=S(Pu)$ both
sides are $Pu$.  Thus $K$ and $W$ are provably equal on the whole domain,
even if they are distinct descriptions.  This equality does not appeal to
standard-model agreement or to uniqueness of an arbitrary recurrence.

\subsection{Four elementary uses of equation induction}

Choose descriptions for addition and truncated subtraction, with recursion
on the second argument:
\begin{align}
 x+0&=x, &x+Sy&=S(x+y),\label{eq:eic-addition}\\
 d(x,0)&=x, &d(x,Sy)&=P(d(x,y)).\label{eq:eic-subtraction}
\end{align}
Using these equations and \eqref{eq:eic-induction}, we now derive
\begin{align}
 0+y&=y,\label{eq:eic-left-zero}\\
 (Sx)+y&=S(x+y),\label{eq:eic-left-successor}\\
 d(Sx,Sy)&=d(x,y),\label{eq:eic-cancel-successors}\\
 d(Sa,a)&=1.\label{eq:eic-diagonal}
\end{align}
Each induction below has a single equation as its induction formula.

For \eqref{eq:eic-left-zero}, induct on $y$.  The base is $0+0=0$.
The step is $0+Sy=S(0+y)=Sy$.

For \eqref{eq:eic-left-successor}, hold $x$ fixed and induct on $y$.
At zero both sides are $Sx$.  At a successor the induction hypothesis
gives
\[
 (Sx)+Sy=S\bigl((Sx)+y\bigr)
       =S\bigl(S(x+y)\bigr)=S(x+Sy).
\]
Combining \eqref{eq:eic-left-successor} with the defining equation for
addition gives the useful rearrangement
\begin{equation}
 (Sy)+r=y+Sr.
 \label{eq:eic-shift-successor}
\end{equation}
Indeed, both sides are $S(y+r)$.

For \eqref{eq:eic-cancel-successors}, hold $x$ fixed and induct on $y$.
The base computes as
\[
 d(Sx,S0)=P(d(Sx,0))=P(Sx)=x=d(x,0).
\]
Applying $P$ to the induction equation gives the step:
\[
 d(Sx,S(Sy))=P(d(Sx,Sy))=P(d(x,y))=d(x,Sy).
\]

Finally, induct on $a$ for \eqref{eq:eic-diagonal}.  The base is
$d(S0,0)=S0=1$.  At a successor,
\[
 d(S(Sa),Sa)=d(Sa,a)=1,
\]
using \eqref{eq:eic-cancel-successors} and the induction hypothesis.

\subsection{An equation invariant that reflects actual equality}

The remaining arithmetic ingredient is the following equation, valid for
all $x,y$:
\begin{equation}
 C\bigl(\operatorname{sg}(d(x,y)),\ x,\ y+d(x,y)\bigr)=x.
 \label{eq:eic-residual-invariant}
\end{equation}
Its meaning is that a nonzero subtraction residual reconstructs the
original argument by addition.  The selector expresses this assertion as
one literal equation, so equation induction applies directly.

Hold $x$ fixed and induct on $y$.  At $y=0$, subtraction gives residual
$x$, and \eqref{eq:eic-left-zero} gives $0+x=x$.  The desired equation
is consequently $C(\operatorname{sg}(x),x,x)=x$, which is
\eqref{eq:eic-selector-idempotence}.

For the step, assume \eqref{eq:eic-residual-invariant} at $y$ and
abbreviate
\[
 r=d(x,y),\qquad r'=Pr=d(x,Sy).
\]
These are local term abbreviations, not additional function symbols.
Split coverage for $r'$.

If $r'=0$, then $\operatorname{sg}(r')=0$, and the successor instance
of \eqref{eq:eic-residual-invariant} reduces immediately to $x=x$.

In the other branch, $r'=S(Pr')$, so constructor separation gives
$r'\ne0$.  Split coverage for $r$.  Its zero branch would imply
$r'=P0=0$, a contradiction, from which the desired step equation follows.
In its successor branch,
\[
 r=S(Pr)=Sr',\qquad \operatorname{sg}(r)=1.
\]
The induction equation therefore computes to $y+r=x$.  Meanwhile,
\eqref{eq:eic-shift-successor} gives
\[
 (Sy)+r'=y+Sr'=y+r=x.
\]
Since $r'$ is a successor, $\operatorname{sg}(r')=1$, so this is
exactly the branch selected in the successor instance of
\eqref{eq:eic-residual-invariant}.  This completes the equation induction.
All case distinctions came from coverage; no induction on an implication
or a disjunction was used.

\paragraph{Reflection at residual $1$.}
We have, for arbitrary $a,b$,
\begin{equation}
 d(Sa,b)=1\quad\longleftrightarrow\quad a=b.
 \label{eq:eic-residual-reflection}
\end{equation}
For the forward direction, substitute $x=Sa$, $y=b$ into
\eqref{eq:eic-residual-invariant}.  If $d(Sa,b)=1$, the sign is $1$
and the selector returns its second branch, giving
\[
 b+1=Sa.
\]
The defining equation for addition gives $b+1=S(b+0)=Sb$.  Thus
$Sb=Sa$, and applying $P$ gives $b=a$, hence $a=b$.
Conversely, $a=b$ and \eqref{eq:eic-diagonal} give
$d(Sa,b)=d(Sa,a)=1$.

Define the equality characteristic by the fixed description
\begin{equation}
 e(a,b)=\operatorname{sg}\bigl(K(d(Sa,b))\bigr).
 \label{eq:eic-equality-characteristic}
\end{equation}
The sign, countdown, and residual reflection lemmas give
\begin{align}
 e(a,b)=0&\quad\longleftrightarrow\quad a=b,
       \label{eq:eic-equality-reflection}\\
 e(a,b)=0&\ \lor\ e(a,b)=1.
       \label{eq:eic-equality-range}
\end{align}
In detail, $e(a,b)=0$ is equivalent to $K(d(Sa,b))=0$ by
\eqref{eq:eic-sign-zero}, which is equivalent to $d(Sa,b)=1$ by
\eqref{eq:eic-countdown-zero}, and then to $a=b$ by
\eqref{eq:eic-residual-reflection}.  The range assertion is
\eqref{eq:eic-sign-range} applied to $K(d(Sa,b))$.

The reflection property concerns the logic's actual equality throughout
the domain.  It was not assumed as part of a numerical comparison symbol,
and its proof did not restrict $a$ or $b$ to standard numerals.
It also yields decidable equality.  In the $e(a,b)=0$ branch,
\eqref{eq:eic-equality-reflection} gives $a=b$.  In the $e(a,b)=1$
branch, an assumption $a=b$ would give $e(a,b)=0$ and therefore $0=1$,
contradicting separation.  Hence
\begin{equation}
 \forall a\,\forall b\,\bigl(a=b\lor(a=b\to\bot)\bigr)
 \label{eq:eic-equality-decidable}
\end{equation}
is a theorem, not an extra hypothesis.

\subsection{Compilation of all quantifier-free formulas}

For each fixed quantifier-free formula $A(\vec p)$, construct a term
$b_A(\vec p)$ by external recursion on its finite syntax.  Take equality,
$\bot$, $\land$, $\lor$, and $\to$ as the primitive formula constructors;
negation is implication into $\bot$.  Set
\begin{align}
 b_{s=t}&=e(s,t), & b_{\bot}&=1,\label{eq:eic-boolean-atoms}\\
 b_{A\land D}&=C(b_A,b_D,1),\label{eq:eic-boolean-and}\\
 b_{A\lor D}&=C(b_A,0,b_D),\label{eq:eic-boolean-or}\\
 b_{A\to D}&=C(b_A,b_D,0).\label{eq:eic-boolean-implies}
\end{align}
Every term here is a finite composition of fixed primitive-recursive
descriptions; if desired, projections and composition give it a single
description symbol.  No internal quantification over formulas or programs
is involved.

An external induction on the syntax of $A$ proves, within the arithmetic
theory,
\begin{equation}
 \bigl(b_A=0\lor b_A=1\bigr)
 \quad\text{and}\quad
 \bigl(b_A=0\longleftrightarrow A\bigr).
 \label{eq:eic-boolean-correctness}
\end{equation}
Here are the full connective checks.  In every case $b_A=1$ implies
$\neg A$: if $A$ held, the already established implication
$A\to b_A=0$ would contradict separation.

\paragraph{Atomic formulas and falsity.}
For an equation, both claims are
\eqref{eq:eic-equality-reflection}--\eqref{eq:eic-equality-range}.
For falsity the range assertion is $1=0\lor1=1$, with its right disjunct
proved by reflexivity.  Also $1=0\to\bot$ follows from separation and
symmetry, while $\bot\to1=0$ follows by explosion.  Thus $1=0$ is
equivalent to $\bot$, as required.

\paragraph{Conjunction.}
Split $b_A=0\lor b_A=1$.  In the zero branch, $A$ holds and
$b_{A\land D}=b_D$.  Its range follows from the range of $b_D$,
and its being zero is equivalent to $D$, hence to $A\land D$ under
the known $A$.  In the one branch, $\neg A$ holds and
$b_{A\land D}=1$.  Its range is immediate.  If it were zero,
separation would give $\bot$, hence $A\land D$; conversely an
assumption $A\land D$ contradicts $\neg A$, and explosion gives
$b_{A\land D}=0$.

\paragraph{Disjunction.}
In the branch $b_A=0$, the formula $A$ holds and
$b_{A\lor D}=0$.  The range is immediate, and both $b_{A\lor D}=0$
and $A\lor D$ hold, proving their equivalence.  In the branch
$b_A=1$, we have $\neg A$ and $b_{A\lor D}=b_D$.  Its range follows
from the range of $b_D$.  If its value is zero, then $D$ and hence
$A\lor D$.  Conversely, from $A\lor D$, the $A$ alternative
contradicts $\neg A$ and yields $b_D=0$ by explosion, while the
$D$ alternative yields $b_D=0$ directly.  This proves the required
equivalence in both alternatives.

\paragraph{Implication.}
In the branch $b_A=0$, the formula $A$ holds and
$b_{A\to D}=b_D$.  Its range follows from the range of $b_D$.
If its value is zero, then $D$, which gives $A\to D$.  Conversely,
$A\to D$ and the known $A$ give $D$, hence $b_D=0$.
In the branch $b_A=1$, we have $\neg A$ and
$b_{A\to D}=0$.  Its range is immediate.  To prove $A\to D$,
assume $A$; this contradicts $\neg A$, and explosion gives $D$.
Thus both $b_{A\to D}=0$ and $A\to D$ hold in this branch, and
their equivalence follows.  This argument uses the proved two-value
range of the finite characteristic.  It does not use excluded middle
for arbitrary formulas or change the surrounding intuitionistic logic.

This completes the external syntactic induction proving
\eqref{eq:eic-boolean-correctness} for every quantifier-free formula.

\subsection{Recovery of induction and the full theory}

Let $A(n,\vec p)$ be any fixed quantifier-free formula.  Suppose an
arbitrary retained context $\Gamma$ proves its base and its induction
step:
\[
 \Gamma\vdash A(0,\vec p),
 \qquad
 \Gamma,A(z,\vec p)\vdash A(Sz,\vec p),
\]
with the usual freshness condition on $z$.  By
\eqref{eq:eic-boolean-correctness}, the base gives
$\Gamma\vdash b_A(0,\vec p)=0$.  Under the additional assumption
$b_A(z,\vec p)=0$, the same equivalence gives $A(z,\vec p)$,
the given step gives $A(Sz,\vec p)$, and the equivalence gives
$b_A(Sz,\vec p)=0$.  Apply literal-equation induction to
$b_A(n,\vec p)=0$ and then apply its reflection implication.
The result is $\Gamma\vdash A(t,\vec p)$ for every target term $t$.
Thus the theory proves the full contextual quantifier-free induction
rule, including in arbitrary quantified contexts.

\paragraph{Completion theorem.}
In the common language of fixed primitive-recursive descriptions, the
theory
\[
 \mathsf{PRE}+\mathsf{Cov}+\mathsf{Sep}+\mathsf{EI}
\]
has exactly the same sequents as ordinary first-order intuitionistic
primitive recursive arithmetic with all quantifier-free induction
instances.  The forward simulation follows from the recovered
quantifier-free induction, successor injectivity, and successor separation;
all primitive-recursive defining equations and all logical rules are
shared.  For the reverse simulation, literal-equation induction is an
instance of quantifier-free induction, and successor separation gives
$\mathsf{Sep}$.  Quantifier-free induction proves $\mathsf{Cov}$:
at zero choose its left disjunct, and at a successor use
$Sx=S(P(Sx))$, which follows from $P(Sx)=x$, for its right disjunct.

Both interpretations therefore use the full domain, defined by $x=x$,
and preserve every function symbol, equality, logical connective, and
quantifier.  Both composites have the identity isomorphism, with graph
$x=y$.  The theories are consequently identity-domain bi-interpretable
and equiconsistent.  This completion theorem is unconditional once the
literal-equation induction rule has been obtained; it has no dependency
on a classical successor-uniqueness reduction.

% End included source


% Begin included source: reference/primitive_recursion.tex
\section{Pairing, primitive recursion, and internal numerical composites}
\label{sec:iterator-coding}
This appendix is adapted with attribution from the progress-closure
arithmetic development~\cite{progress}. Its prerequisites have now been
proved: scalar predecessor, arithmetic reflection, and full quantifier-free
induction. Pairing is a derived proof tool; it is not part of any checker.

\subsection{Pairing from a bounded square-root search}

Once QFI is available, a familiar square-based injection gives a
particularly direct compiler. Write $a\le b$ for $d(a,b)=0$.
The usual linear-order and semiring identities follow by fixed
quantifier-free inductions from the successor, addition, and monus
equations. Multiplication itself is the scalar term
\[
                        a\cdot b=I(b,0;u.u+a).
\]
Only primitive-recursive arithmetic is involved in the following
bounded root search; there is no minimization over an unbounded domain.

For a fixed parameter $n$, put
\[
 U_n(v)=C(d((Sv)^2,n),Sv,v),\qquad
 R_n(i)=I(i,0;v.U_n(v)),\qquad q(n)=R_n(n).
\]
The update increments $v$ exactly when $(v+1)^2\le n$, and
otherwise retains it. QFI in $i$, with $n$ fixed, proves
\[
 R_n(i)\le i,\qquad R_n(i)^2\le n,\qquad
 R_n(i)=i\ \vee\ n<(S(R_n(i)))^2.              \tag{Root invariant}
\]
The base has $R_n(0)=0$. If the update increments, its test
contradicts the strict-upper-bound alternative in the third invariant.
Thus the old state equals $i$, and the new state equals $Si$;
its square bound is the increment test. If the update does not increment, the
strict upper square bound holds and the state remains fixed.
An already established strict upper square bound forces every
subsequent update to remain fixed. These are exactly the cases
needed for the three displayed invariants.

At $i=n$, either the strict upper square bound already holds or
$R_n(n)=n$, in which case $n<(n+1)^2$ supplies it. Thus
\[
                        q(n)^2\le n<(q(n)+1)^2.           \tag{Root}
\]
Monotonicity of squares shows uniqueness of a number satisfying
Root. All formulas just used are quantifier free in the scalar
language, with fixed $n$ during the induction.

Now define
\[
 \pi(a,b)=(a+b)^2+a,\qquad
 L(n)=d(n,q(n)^2),\qquad H(n)=d(q(n),L(n)).
\]
For $k=a+b$, one has $k^2\le k^2+a<(k+1)^2$ since $a\le k$.
Root uniqueness gives $q(\pi(a,b))=a+b$, and therefore
\[
                    L(\pi(a,b))=a,\qquad H(\pi(a,b))=b.
\]
Surjectivity of $\pi$ is unnecessary. The original numerical domain
is never replaced by the range of this coding; pairing is used only
to store the finite state of a fixed primitive recursion. This
standard square pairing is a convenient reconstruction here, not
a claimed new arithmetic basis.

\subsection{Compiling each primitive recursion}

Suppose the fixed target description is
\[
 f(0,\vec p)=g(\vec p),\qquad
 f(Sn,\vec p)=h(n,f(n,\vec p),\vec p).
\]
Assume by external structural induction that the earlier $g,h$
already have scalar translations $g^*,h^*$. Define
\[
 W(u,\vec p)=\pi(SL(u),h^*(L(u),H(u),\vec p)),
\]
\[
 V(n,\vec p)=I(n,\pi(0,g^*(\vec p));u.W(u,\vec p)),
 \qquad f^*(n,\vec p)=H(V(n,\vec p)).
\]
The update has only its state variable bound, with every $p_i$
fixed. Equation induction proves $L(V(n,\vec p))=n$.
The two projection equations and scalar computation then prove
both defining equations for $f^*$ internally. Initial functions
and composition translate directly. This completes the translation
of every fixed finite PR description without adding a universal
evaluator.

\subsection{Why both numerical round trips are internal}

Two fixed scalar iterators with the same seed and provably equal
updates are internally equal at all counters: use literal-equation
induction comparing their values. Their base agrees, and their
successor equations plus congruence prove the step. The same
uniqueness argument works for ordinary primitive recursion.

For an ordinary PR description, induct externally on its finite
construction. The double translation satisfies its original
defining equations, so internal equation induction identifies it
with the original operation. For a native scalar description,
earlier translated updates agree by the external induction
hypothesis; scalar uniqueness identifies the double-translated
iterator. These proofs hold at arbitrary internal counters.

Thus the domain relation $x=y$ is provably a total bijection
preserving every fixed numerical operation in both composites.
The formula constructions and round trips in
Section~\ref{sec:choice-introduction} give the corresponding formula
composites. Both interpretations preserve falsity, taking native
$0=1$ to an ordinary contradiction and ordinary $\bot$ to native
$0=1$. Finite proof translations therefore establish equiconsistency.

% End included source


\clearpage
\section*{Reference proof for the original arity codec}
\addcontentsline{toc}{section}{Reference proof for the original arity codec}
The following section is reproduced unchanged from the arity-calculus
archive~\cite{arity}. It concerns that archive's original computation
packet and original verdict driver. The exact conversion in
Section~\ref{sec:composition}, rather than an identification of different
programs, transfers its soundness argument to the selected checker.
The reference's statements about its own measured profiles are not the
measurements of the present program.

% Begin included source: reference/arity_codec_proof.tex
\section{A uniform arity encoding and its checking theorem}
\label{sec:arity-checker}

The arithmetic and object logic remain the scalar-iteration Eq/$\D$
calculus specified above.  The change is in its finite certificate
representation.  A paired formula is unary, an equation is binary, and
an internal closure is ternary:
\[
 \D_{x,y}B\longleftrightarrow[B],\qquad
 s=t\longleftrightarrow[s,t],\qquad
 \Q_yB\longleftrightarrow[B,B,B].
 \tag{Arity encoding}\label{eq:arity-encoding}
\]
The three copies in a closure are representation redundancy: only the
third has mathematical content.  Their benefit is that tripling a unary
major constructs exactly its closure.  Tripling any other nonempty
formula array produces an array of arity at least six, which will denote
truth.  This makes the other major cases neutral without a separate
test.

The previous generated-assumption and constructed-computation interfaces
are inherited from the corpus~\cite{generated,constructed}.  The new
encoding makes a single arithmetic expression, $5-\operatorname{arity}$,
perform the necessary depth adjustments.  Its soundness depends on an
explicit interpretation of surplus syntax.  It would not be sound to
drop arity checks while continuing to treat every surplus array as an
unconstrained new connective.

\subsection{The complete sorted interpretation of arrays}
\label{subsec:sorted-arrays}

Write $\top$ for the fixed equation $0=0$.  A protected numerical slot
has width two.  The numerical root has depth three: the coordinates
$0,1,2$ all denote zero, and under $p$ protected slots the depth is
$3+2p$.  Coordinates $3+2j$ and $4+2j$ name slot $j$, counted from
the innermost slot, for $0\leq j<p$.  Formula depth is $2p$ under the
same $p$ slots.  Negative coordinates admitted by the current global
floor name rigid numerical parameters.

The decoder is sorted by the position in which a tree is used.  Let
$\rho$ be its ordered environment of protected numerical slots.
The notation $(a,\rho)$ means prepend a slot of value $a$.
A superscript $\sharp$ denotes a child decoded in the indicated
environment.

\begin{center}
\begin{tabular}{cll}
\toprule
Numerical arity & Meaning & Relevant child environment\\
\midrule
$0$ & $0$ & No child\\
$1$: $[b]$ & $Sb^\sharp$ & $(0,0,\rho)$\\
$2$ & $0$ & Annotations ignored\\
$3$: $[n,a,h]$ & $I(n^\sharp,a^\sharp;v.h^\sharp)$ & $(0,\rho)$ for $n,a$; $(v,\rho)$ for $h$\\
$\geq4$ & $0$ & All children ignored\\
\bottomrule
\end{tabular}
\end{center}

\begin{center}
\begin{tabular}{cll}
\toprule
Formula arity & Meaning & Relevant child environment\\
\midrule
$0$ & $\top$ & No child\\
$1$: $[B]$ & $\D_{x,y}B^\sharp$ & $(y,x,\rho)$\\
$2$: $[s,t]$ & $s^\sharp=t^\sharp$ & $\rho$ numerically\\
$3$: $[P,U,B]$ & $\Q_yB^\sharp$ & $(y,\rho)$ for $B$\\
$\geq4$ & $\top$ & All children ignored\\
\bottomrule
\end{tabular}
\end{center}

The first two entries of a ternary formula are annotations.  A canonical
closure copies its body into both of them, as in
\eqref{eq:arity-encoding}.  Successfully admitted scalar formula padding
also denotes $\top$.  An internal closure has the ordinary meaning
$B(0)\to\forall yB(y)$; it is not an equation and is not eligible for
compound induction.

Canonical encoded object expressions use arrays and integer coordinates,
respect these slots, and contain no high-arity node.  The checker itself
still generates high-arity truth entries in administrative contexts,
even during a canonical proof.  In particular, constructing a
successor beneath positive bound coordinates requires passing through
two dummy zero slots.  Constructing an iterator requires one new slot:
zero in its counter and seed, and its genuine state variable in the
update.  Rigid negative parameter names are unaffected by these shifts.

The additional raw shapes are certificate aliases with the displayed
meanings.  They are not independent predicate or function symbols added
to the object theory.  Our exact completeness statement concerns the
canonical native language.  We do not require every possible raw spelling
of a semantic alias to be interchangeable under the same raw certificate.

\subsection{One depth calculation replaces the lookup table}

For every array the opener recursively adds $5-\operatorname{arity}$
to its depth.  The meaningful cases are
\[
\begin{array}{c|ccc}
 \operatorname{arity}&1&2&3\\\hline
 \text{depth increment}&4&3&2.
\end{array}
\]
Unary $\D$ adds two numerical binders; unary successor adds two dummy
zero slots.  Binary equality changes formula depth $2p$ to numerical
depth $3+2p$.  Ternary $\Q$ adds one numerical binder; ternary iteration
adds one state slot, or one dummy slot in its counter and initializer.
The same calculation therefore covers all relevant constructors.

Binary numerical annotations denote zero, so their parity-changing
traversal has no mathematical effect.  Empty arrays have no children.
At arity at least four, the entire node already denotes a constant.
Even a zero or negative depth increment below such a node cannot expose
one of its ignored children as a meaningful variable.  Traversal may
still reject an alias; it cannot change its stipulated interpretation.

This is a semantic assertion, not raw idempotence of the opener.
For example, at floor $-1$, opening a six-element array of zeros with
image $-1$ can yield six copies of $-1$, while a later no-image opening
of that result fails at a negative child depth.  Both formula arrays
denote $\top$.  The proof below only concerns successful calls at meaningful
sorted positions.  Canonical input object expressions contain no surplus
arrays; the administrative exports that generate them are justified
separately as truth entries.

\subsection{Opening and scope}
\label{subsec:opening-lemma}

Write $i(a,t,d)$ for the executable opener at floor $k$.  It preserves a
scalar when $k\leq a<d$.  At a numerical position, every other relevant
scalar is the distinguished substitution hole; the supplied image is
checked at numerical root depth three before insertion.  At a formula
position scalar padding is truth and carries no numerical hole.
Missing images at used numerical holes reject.

\begin{lemma}[Semantic opening]
\label{lem:semantic-opening}
For every successful call at a meaningful formula root of depth $2p$
or numerical root of depth $3+2p$, with $p\geq0$, decoding its result
is capture-free filling of the distinguished hole in the decoded input
template.  These depths protect $p$ numerical slots; admitted rigid
parameters are preserved.  If an image is semantically unused, zero
may be chosen as its mathematical representative.  Recursive calls
at negative depths inside ignored annotations need no independent
sorted interpretation.
\end{lemma}

\begin{proof}
Induct on the finite raw tree.  At a relevant scalar, the interval
comparison is exactly the protected-coordinate test.  At a meaningful
unary, binary, or ternary node, the preceding tables supply precisely
the required additional slots.  At numerical binary annotations,
high-arity nodes, and ignored closure fields, the relevant mathematical
meaning is constant or absent, so the conclusion is immediate
regardless of traversal inside the annotation.

An inserted image is checked at its own numerical root.  Its free
parameters are rigid negative names, and it has no unbound positive
coordinate there.  Insertion beneath protected slots therefore captures
none of its parameters; its internal binders retain their relative
indices.  This proves the inductive step for substitution as well as
for mere traversal.
\end{proof}

A successful no-image check consequently preserves the interpretation
of its checked closed expression.  Raw structural equality of checked
outputs implies equality of their decoded meanings.  The checker does
not infer equality merely from a semantic comparison of two different
programs or numerical terms.

The global floor decreases by one at each non-hypothesis continuation.
If its old value is $K$, the integer $K-1$ is the designated fresh
eigenparameter.  With integer coordinates it is the only new name.
If fractional coordinates are accepted, other names in $[K-1,K)$ may
also occur in the finite certificate.  These are additional fresh
parameters.  They must not be silently equated to the designated
integer eigenparameter.

Previously checked relevant goals and retained hypotheses mention none
of these new names: their earlier no-image checks would reject a
relevant numerical occurrence below the old floor.  Coordinates inside
ignored annotations do not contribute free variables to the decoded
formula.  This is the freshness invariant used in every rule below.

\subsection{A computation hypothesis only needs a zero instance}
\label{subsec:zero-supply}

The shared continuation can open a generated arithmetic equation at a
fresh parameter.  That equation need not be a theorem uniformly in the
new parameter.  The following admissible rule is sufficient.

\begin{lemma}[Supply with a known zero instance]
\label{lem:zero-supply}
Suppose $\vec a$ is fresh for $\Gamma,F$, the calculus proves
$H(\vec0)$, and $\Gamma,H(\vec a)\vdash F$.  Then
$\Gamma\vdash F$.
\end{lemma}

\begin{proof}
Substitute zero for every member of $\vec a$ in the second derivation.
Freshness leaves $\Gamma,F$ unchanged.  Cut in the derivation of
$H(\vec0)$.  Both substitution and cut are admissible in the native
calculus.  This is a finite proof transformation and does not use a
consistency assumption.
\end{proof}

All names newly admitted when the floor changes belong to $\vec a$;
for noninteger aliases this includes the finite collection in
$[K-1,K)$ that actually occurs.  On decoding, their simultaneous zero
substitution agrees with opening the generated expression at zero using
the old floor.  Previously admitted parameters are preserved.  New
rigid names and unprotected placeholders both become zero.  The opener
may keep different raw zero spellings, but their decoded values agree.

The constructed-computation packet supplies a triple $s=[n,b,r]$ and a
mode $u$.  In successor mode its generated equation has raw form
\[
 \bigl[\,[[n],b,r],\ i(r,s,3)\,\bigr],
 \tag{Step supply}\label{eq:step-supply}
\]
while in zero mode it has form
\[
                    \bigl[\,[u,b,r],\ b\,\bigr].
 \tag{Zero supply}\label{eq:base-supply}
\]
Here successor mode means that $u$ is truthy; the zero mode has falsy
$u$.  These expressions describe actual arrays, before their final
opening.  In \eqref{eq:step-supply} the counter is constructed as a
successor, rather than accepted on an unsupported assertion that it is
one.  This construction is inherited from the corresponding draft
interface~\cite{constructed}.

\begin{lemma}[Sound constructed computation]
\label{lem:constructed-supply}
Every successful arithmetic supply has a decoded zero instance
derivable from the two scalar-iteration computation schemas.
\end{lemma}

\begin{proof}
First apply the simultaneous zero substitution described above.  In
zero mode every successful falsy numerical counter denotes zero, so
\eqref{eq:base-supply} is an $I0$ instance.  In the exact-list profile
the relevant accepted aliases are an empty array, numeric zero, false,
and numerically equal zero values.  Null, strings, and dictionaries
cannot enter a meaningful numerical position through an unchecked
coercion.

In successor mode the left counter is an actual successor.  If the
update uses its state in a semantically relevant position, opening it
invokes the no-image check of the recursive source $s$ at numerical
root three.  This checks the entire source, including its update.
Only coordinates three and four are aliases for its single state slot.
A genuinely escaping coordinate in that update is revisited during
the recursive-image check and rejects.  Its counter and initializer
are likewise closed with their prescribed dummy zero slot.

Inserting the extra dummy slots of the displayed successor does not
alter such a checked predecessor: extra unbound slots have value zero,
and actual free parameters are rigid.  Lemma~\ref{lem:semantic-opening}
therefore identifies the right side of \eqref{eq:step-supply} with the
update applied to exactly the recursive predecessor value occurring
on its left.  This is the $IS$ schema.

If the update has no semantically relevant state occurrence, its
decoded value is constant in that variable.  Its zero-normalized value
is the same in the source update and on the right, independently of
the unused recursive image; this is again an $IS$ instance.  Traversal
inside an ignored annotation may still inspect the image and reject.
That removes a certificate and does not invalidate the argument for a
successful one.

Finally, successful three-field source extraction cannot hide another
JSON operator.  A scalar is not iterable.  A three-character string,
or a three-key JSON dictionary, necessarily puts a nonempty string in
a relevant source field, which the inspected opener rejects.  The
successful source is therefore the stated array triple.
\end{proof}

The recursive-image check in the state-dependent case is essential.
Merely opening a final displayed equation at zero would not establish
that its two state interpretations agree.  Combining
Lemma~\ref{lem:constructed-supply} with Lemma~\ref{lem:zero-supply}
justifies both actual image choices of the common continuation, including
the branch that retains a fresh-parameter hypothesis.

\subsection{Audit of the logical continuations}
\label{subsec:logical-audit}

\paragraph{Hypothesis closure.}
The initial opener check and membership comparison select the same
decoded context formula.  Additional aliases can make a spelling less
convenient, but cannot identify an equation with a false denotation
solely by semantic evaluation.

\paragraph{Equality transport.}
The major must match $[s,u]$, and opening the proposed formula template
at $u$ must match the target.  The continuation checks its opening at
$s$.  Lemma~\ref{lem:semantic-opening} gives ordinary transport along
the decoded equality.  The argument applies to compound formulas,
internal clauses, and truth aliases as well as equations.

\paragraph{Closure specialization.}
The closure pattern is the triplicate $[t,t,t]$.  The guard forces the
old instance to have numerical image zero; the new instance may have
an arbitrary numerical image.  This is exactly
$B(0)\to\forall yB(y)$.  The zero guard is indispensable.  Without
it, a closure with false zero instance and true successor instances
would turn a proved successor instance into falsity.

\paragraph{Paired elimination.}
The export is always three concatenated copies of the major.  A unary
major $[B]$ therefore exports $[B,B,B]$.  All three copies are opened
identically; the third copy has precisely the witness instance required
by paired elimination.  The newly allocated witness is fresh for the
retained context and conclusion.

An equation major instead exports arity six; a closure major exports
arity nine; a larger array exports still larger arity.  These are all
truth aliases.  Thus failed elimination endpoint claims have only the
effect of adjoining a theorem.  A numeric scalar major, if accepted as
formula padding, exports scalar padding after multiplication and also
denotes truth.  Unsupported types reject.

A ternary formula with unequal annotations still denotes the closure
of its third child.  It fails the triplicate specialization pattern
unless the three fields match.  This loses a possible raw spelling,
not a mathematical soundness condition.  Canonical compilation uses
triplicate closures throughout.

\paragraph{Paired introduction.}
The call \texttt{i(*a,t,2)} uses the function's argument count to enforce
the relevant outer arity.  For a unary goal $a=[B]$, it is exactly
\texttt{i(B,t,2)}: the universal slot is protected while the
existential witness is opened.  For arity at least two, there are too
many arguments and the call rejects.  An empty goal can instead make
the legal two-argument call \texttt{i(t,2)}.  Its original target was
already $\top$, so accepting it cannot add a false theorem.

For the unary goal, the subsequent successor opening gives the intended
premise $B(t,Sz)$.  The local hypothesis is $B(t,0)$ in case mode and
$B(t,z)$ in progress mode.  The freshness invariant supplies the
usual eigenvariable conditions.

\paragraph{Induction remains restricted.}
Among meaningful nonconstant matrix constructors, only binary equality
has even arity.  Unary paired formulas and ternary closures have odd
arity and receive case introduction.  Empty or surplus arrays that
reach the progress branch denote the fixed equation $0=0$, whose
induction is valid.  Scalar padding that cannot provide a length
rejects.  The broader certificate alphabet therefore adds no induction
on a quantified matrix.

\paragraph{Additional fresh parameters.}
An accepted subderivation can contain newly admitted fractional names,
but the older context and target do not depend on them.  Replace these
auxiliary names by zero when closing the rule, retaining its designated
integer eigenvariable where that rule needs one.  The common opening
of an already checked retained goal leaves its decoded meaning
unchanged.  This handles the actual global-floor convention without
pretending that every admissible negative number is an integer.

\subsection{Canonical completeness, including zero witnesses}
\label{subsec:canonical-completeness}

Rename a finite native proof's eigenvariables to fresh negative integers.
Encode its terms with the slot conventions of
Section~\ref{subsec:sorted-arrays}; encode paired formulas as unary and
closures as triplicate.  Hypothesis leaves consume no instruction.
Transport, closure specialization, and paired elimination have their
displayed packets.  Arithmetic computation leaves use
\eqref{eq:base-supply} or \eqref{eq:step-supply} and then close the
target by hypothesis lookup.

One representational issue requires an actual proof transformation.
The packet's truthy witness field selects introduction.  A literal zero
witness therefore cannot simply be inserted into that field.  Define
\[
                 J(t)=I(0,t;v.0).
\]
The computation axiom proves $J(t)=t$, and $J(t)$ has a nonempty array
representation even when its value is zero.  To introduce with a falsy
witness spelling $t$, introduce instead with $J(t)$: transport its
retained assumption to the assumption using $t$, use the original
premise derivation, and transport the resulting successor instance
back to the instance using $J(t)$.  Derived symmetry supplies the
reverse transport.  This is a finite transformation of the premise,
and the paired conclusion itself is unchanged.  If the existential
coordinate is unused, any closed truthy witness works without a
conversion.

The scalar coordinate $1$ is also a semantic zero at numerical root
three, but that observation alone is not a completeness proof.  Closure
specialization and case introduction insert literal raw zero, and
distinct zero aliases need not be interchangeable under a fixed raw
proof spelling.  The computation-backed box gives the required
derivation instead of assuming representational interchangeability.

The mode zero selects progress at an equation.  A different falsy
mode, such as the empty array, selects case.  These are certificate
control tags, not numerical assertions about the witness.  All steps
in the compilation are finite and effective.

\begin{theorem}[Conservativity and completeness of the arity checker]
\label{thm:arity-checker}
Every accepted certificate has its decoded conclusion provable in the
ordinary realization of the native calculus.  For every canonical
native end sequent, the internal checking judgment and the native
calculus have the same consequence relation, in the idealized
unbounded integer-and-array implementation.  In particular, the public
empty-context driver recognizes exactly the canonical closed theorems.
The translation in each direction is effective on finite derivations.
\end{theorem}

\begin{proof}
For soundness, induct on the finite accepting call tree.  Hypothesis
closure uses Lemma~\ref{lem:semantic-opening}; the logical cases are
the audited continuations above.  Arithmetic supply uses
Lemmas~\ref{lem:constructed-supply} and \ref{lem:zero-supply}.
Every discarded annotation is accounted for by its zero or truth
interpretation, and every fresh parameter by the scope invariant.
For completeness, compile a finite native derivation by the preceding
construction, including the boxed-witness transformation.  Its
canonical input expressions require no surplus-arity aliases or fractional
names; generated administrative surplus arrays have exactly the truth
meaning already proved.  This proves both implications without using a finite test
suite as a substitute for induction on derivations.
\end{proof}

The unchanged native calculus is bi-interpretable with full first-order
$\ipra$ by the inherited arithmetic and logical bridge
\cite{generated,corpus}.  In particular,
\[
 (\D_{x,y}B)^\circ
 =\exists x\bigl(B^\circ(x,0)\to\forall yB^\circ(x,y)\bigr).
\]
The whole numerical domain and actual equality are retained.  Its
internally proved numerical round trips and its formula-composite
derivations therefore transfer through
Theorem~\ref{thm:arity-checker}.  We have changed the checked
representation, not replaced bi-interpretability by standard-model
agreement or supplied a new proof of the scalar predecessor theorem.

\subsection{JSON and finite-runtime qualifications}

The primary profile uses an ordinary list-type test.  The smaller byte
profile instead tests multiplication by zero against an empty list.
Among standard JSON runtime values the latter recognizes exactly lists:
numeric values and booleans produce numeric zero, strings produce an
empty string, and null and dictionaries raise rather than becoming
arrays.  This argument assumes ordinary decoded JSON values, not
arbitrary Python objects with overloaded methods.

Finite nonnegative real coordinates have the same interval
interpretation: $[0,3)$ is zero, and $[3+2j,5+2j)$ names slot $j$.
Finite negative coordinates in the admitted range are additional rigid
names.  Equal Python numerical values have the same decoded role,
including Boolean aliases and signed zero.  Nonfinite values admitted
by Python's permissive JSON reader cannot survive a relevant closed
numerical check.  They may be filled as template holes or ignored
inside a constant annotation; neither makes them unchecked parameters.

The alternative opener based on generic iteration has a separate
runtime qualification.  Empty strings and empty dictionaries traverse
as empty arrays; a dictionary whose sole key is empty traverses as a
unary array with an empty child.  Dictionary values are not visited.
Nonempty strings can rely on CPython's recursion limit to reject, as
also occurs in a current corpus comparison profile~\cite{transient}.
The idealized structural-termination claim for the exact-list profile
must not be transferred to that generic profile unchanged.

All executable profiles have ordinary finite memory, stack, and input
limits.  These can reject a sufficiently large otherwise valid proof.
The formal consequence relation is the unbounded finite-certificate
calculus; the delivered program checks certificates within its runtime
resources.  The measurements charge the entire delivered checker,
while the semantic proof and canonical producer are supplied separately.

% End included source


\clearpage
\begin{thebibliography}{99}\raggedright
\bibitem{snapshot}
Vladimir Reshetnikov, repository maintainer.
\emph{ProveIt, current draft collection}, commit
\file{fb2dbd2a33c6b132671a54b029d0bf3cd98c633d}.
\url{https://github.com/VladimirReshetnikov/ProveIt/tree/fb2dbd2a33c6b132671a54b029d0bf3cd98c633d/docs/drafts}.
All archive references below refer to this immutable collection.

\bibitem{uniform}
\emph{Uniform Predecessor Certificates for Intuitionistic Primitive
Recursive Arithmetic}. \file{ipra_uniform_predecessor_2026-11-10.zip},
in~\cite{snapshot}. Source of the predecessor instruction, current
comparison leaders, and attributed complete logical/arithmetic appendices.

\bibitem{arity}
\emph{An Arity Calculus for iPRA}.
\file{ipra_arity_calculus_2026-11-10.zip}, in~\cite{snapshot}.
Source of the unary paired binder, triplicate closure, numerical root
depth three, and $5-\operatorname{arity}$ traversal.

\bibitem{research}
\emph{Intuitionistic primitive recursive arithmetic research report}.
\file{ipra_research_2026-11-10.zip}, in~\cite{snapshot}, especially
\file{theory/finite_parameter_obstruction.tex}, the eventually constant power and its
finite-conjunction generalization. The conjunctive transfer theorem in
this report extends that construction.

\bibitem{affine}
\emph{Affine Depths for a Small iPRA Checker}.
\file{ipra_affine_arity_2026-11-10.zip}, in~\cite{snapshot}.

\bibitem{generated}
Inherited generated-assumption interface, reproduced and attributed in
the current arity-calculus and uniform-predecessor archives,
in~\cite{snapshot}.

\bibitem{constructed}
\emph{Constructing computation sources for iPRA}.
\file{ipra_constructed_computation_2026-11-10.zip}, in~\cite{snapshot}.

\bibitem{transient}
\emph{A one-state recursor for iPRA}.
\file{ipra_transient_seed_2026-11-10.zip}, in~\cite{snapshot}.

\bibitem{corpus}
The pinned draft corpus~\cite{snapshot}; this bibliography key is retained
for the unchanged arity-codec reference section.

\bibitem{erasure}
Inherited scalar-predecessor development, reproduced with attribution
as \file{reference/scalar_predecessor.tex} in~\cite{uniform} and here.

\bibitem{forward}
Inherited forward-left equality-reflection development, reproduced with
attribution as \file{reference/equality_reflection.tex} in~\cite{uniform}
and here.

\bibitem{progress}
Inherited pairing and primitive-recursion compiler, reproduced with
attribution as \file{reference/primitive_recursion.tex} in~\cite{uniform}
and here.

\bibitem{avigad}
Jeremy Avigad. \emph{Proof mining}, lecture slides, 2004, slide 13.
\url{https://www.andrew.cmu.edu/user/avigad/Talks/asl04.pdf}.
Used only for the standard background distinction between first-order
classical, first-order intuitionistic, and quantifier-free presentations.

\bibitem{moschovakis}
Yiannis N. Moschovakis. \emph{Notes on Set Theory and Logic},
Chapter 3, Section 3J, primitive recursive arithmetic.
\url{https://www.math.ucla.edu/~dam/220b.1.11w/220notes3.pdf}.
Used only as a background reference for staged primitive-recursive
descriptions and quantifier-free induction.
\end{thebibliography}
\end{document}
